\documentclass[11pt]{amsart}

\usepackage[margin=1in]{geometry}
\usepackage[T1]{fontenc}
\usepackage{lmodern}
\usepackage{amsmath,amssymb,amsthm,mathtools}
\usepackage{microtype}
\usepackage{booktabs}
\usepackage{url}
\usepackage{array}
\usepackage{enumitem}
\usepackage[dvipsnames]{xcolor}
\usepackage[colorlinks=true,linkcolor=MidnightBlue,citecolor=OliveGreen,urlcolor=MidnightBlue]{hyperref}
\usepackage[capitalize,nameinlink]{cleveref}

\setlist{leftmargin=*,itemsep=2pt,topsep=4pt}
\emergencystretch=2.5em
\setlength{\tabcolsep}{5pt}
\renewcommand{\arraystretch}{1.12}

% ---------------------------------------------------------------- theorems
\theoremstyle{plain}
\newtheorem{theorem}{Theorem}[section]
\newtheorem{lemma}[theorem]{Lemma}
\newtheorem{proposition}[theorem]{Proposition}
\newtheorem{corollary}[theorem]{Corollary}
\newtheorem{fact}[theorem]{Fact}
\newtheorem{importedlemma}[theorem]{Imported Lemma}
\newtheorem{importedtheorem}[theorem]{Imported Theorem}
\newtheorem*{hypA}{Hypothesis $(\mathrm{H}_{\mathrm{A}})$}
\theoremstyle{definition}
\newtheorem{definition}[theorem]{Definition}
\newtheorem{conjecture}[theorem]{Conjecture}
\theoremstyle{remark}
\newtheorem{remark}[theorem]{Remark}

\crefname{importedlemma}{Imported Lemma}{Imported Lemmas}
\crefname{importedtheorem}{Imported Theorem}{Imported Theorems}
\crefname{fact}{Fact}{Facts}
\crefname{conjecture}{Conjecture}{Conjectures}
\crefname{step}{Step}{Steps}
\numberwithin{equation}{section}

% ---------------------------------------------------------------- macros
\newcommand{\HC}{\#\mathrm{HC}}
\newcommand{\HP}{\#\mathrm{HP}}
\newcommand{\bpsi}{\bar\psi}
\newcommand{\Z}{\mathbb{Z}}
\newcommand{\Q}{\mathbb{Q}}
\newcommand{\F}{\mathbb{F}}
\newcommand{\cL}{\mathcal{L}}
\newcommand{\cF}{\mathcal{F}}
\newcommand{\topc}[2]{[#1]\bigl(#2\bigr)}
\newcommand{\lab}[1]{\mathsf{#1}}
\newcommand{\Eint}{E_{\mathrm{int}}}
\newcommand{\mx}{m_{\times}}
\newcommand{\enc}{\operatorname{enc}}
\newcommand{\block}{\operatorname{block}}
\newcommand{\poly}{\operatorname{poly}}
\newcommand{\supp}{\operatorname{supp}}
\newcommand{\cc}{\operatorname{cc}}
\newcommand{\Ind}[1]{[\,#1\,]}
\DeclareUrlCommand\lean{\urlstyle{tt}}
\newcommand{\A}[1][]{\textup{[A\if\relax\detokenize{#1}\relax\else, #1\fi]}}

\begin{document}

\title[Counting Hamiltonian cycles in sparse graphs]{Counting Hamiltonian cycles in sparse graphs\\ in subexponential time,\\ assuming the new $\#$SAT algorithms}

\author{Vaibhav Tulsyan}

\date{October 2026}

\subjclass[2020]{68Q25, 68W40, 05C45, 05C30, 15A75}
\keywords{Hamiltonian cycle, exact counting, Exponential Time Hypothesis, Grassmann algebra, matrix-tree theorem, Cut\&Count, tensor networks}

\begin{abstract}
In October 2026 two manuscripts claimed refutations of the Exponential Time Hypothesis by
deterministic algorithms for $\#k$-SAT, each with a Lean formalisation of its top-level
statement. Neither result applies directly to connectivity problems such as Hamiltonicity,
because no linear-size CNF encoding of ``this edge set is a single cycle'' is known. This
paper derives a corollary of the Anthropic manuscript: assuming three of its internal results (a
network-to-clique reduction, an algorithm for cliques of all levels, and the existence of its
``pieces''), the number of Hamiltonian cycles of a graph with $n$ vertices and at most $Dn$
edges can be computed exactly and deterministically in time $O(2^{\delta n})$ for every
$\delta>0$, and in time $2^{O(\max(m,n)/\sqrt{\log n})}$ for a graph with $m$ edges. A
rank-encoded variant of the construction gives $2^{O(n\log\bar d/\sqrt{\log n})}$ for average
degree $\bar d$. Deterministic decision, search and Hamiltonian-path counting follow, and so do
exact sparse TSP with polynomially bounded weights, exact counting of directed Hamiltonian cycles,
and ATSP. The result is not a
standalone breakthrough; it stands or falls with the claimed refutation. The ingredients are
known: the folklore matrix-tree identity on $2$-factors, the fermionic (Berezin) matrix-tree
theorem of Caracciolo, Jacobsen, Saleur, Sokal and Sportiello, and the blob partitions of
Cut\&Count. What is new is the observation that connectivity can be encoded with $O(1)$ bits
per edge through a Grassmann determinant whose only non-local part is a sign, and that this
sign costs only $O(g^3)$ bits for $g$ blobs, so the whole construction is a network in the sense
of the manuscript, whose value is $n\cdot\HC(G)$. Every theorem we claim is formally verified in Lean, end to end, on top of the manuscript's own
Lean development: the three counting theorems, the decision, search and path corollaries, sparse
TSP, directed Hamiltonian cycles and ATSP, using only the standard axioms, with fresh kernel replays. We also record weaker results
that rest only on the machine-checked top-level theorems: the parity of $\HC(G)$ for every
graph in time $2^{O((m+n)/\sqrt{\log(m+n)})}$ (machine-checked in Lean, with the hypotheses
discharged by the labs' own Lean proofs), and one-sided randomized decision. We describe what
is proved, what is checked numerically, what is formalised, what is assumed, and what the
limitations are.
\end{abstract}

\maketitle

\setcounter{tocdepth}{1}
{\small\tableofcontents}

% =====================================================================
\section{Introduction}\label{sec:intro}
% =====================================================================

\subsection{Background}
The Exponential Time Hypothesis (ETH) of Impagliazzo and Paturi~\cite{ImpagliazzoPaturi01}
asserts that there is $\delta>0$ such that $3$-SAT on $n$ variables cannot be solved in time
$O(2^{\delta n})$. Together with the sparsification lemma~\cite{IPZ01} it underlies a large body
of conditional lower bounds~\cite{LMS11,CyganEtAl15}. Among them is the statement that, unless ETH
fails, Hamiltonian Cycle has no $2^{o(n+m)}$-time algorithm on graphs with $n$ vertices and $m$
edges~\cite[Ch.~14]{CyganEtAl15} (see also \cite[Thm.~3.3]{LMS11}, stated there as the
absence of a $2^{o(n)}$-time algorithm for Hamiltonian Path).

In October 2026 two manuscripts appeared, by OpenAI~\cite{OpenAI26} and by
Anthropic~\cite{Anthropic26}, each claiming a deterministic algorithm that counts the satisfying
assignments of a $k$-CNF in subexponential time, and hence a refutation of ETH. Both are
accompanied by Lean~4 formalisations of their top-level statements. We stress that these are
\emph{claimed} results that are, at the time of writing, still under review by the community.
We write \A{} for the Anthropic manuscript~\cite{Anthropic26}. Its Main Theorem gives one
deterministic algorithm that counts the models of a $k$-CNF with $m$ clauses on $n$ variables in
time
\[
  2^{O(\min\{km/\sqrt{\log m},\;k^2 n\log\log n/\sqrt{\log n}\})}
\]
with an absolute constant in the $O$. The bound of~\cite{OpenAI26} that is machine-checked is
$2^{O_k(\rho(n))}$ with $\rho(N) = N(\log\log N)^2/\sqrt{\log N}$.

A refutation of ETH does not by itself yield fast algorithms for connectivity problems. ETH
lower bounds for Hamiltonicity come from linear-size reductions \emph{from} $3$-SAT \emph{to}
Hamiltonian Cycle; an algorithm needs the reverse direction, an encoding of Hamiltonicity as a
counting problem on $O(n)$ Boolean variables with short clauses. Connectivity is a global
property and no such encoding is known. Accordingly, the consequences listed by both labs
(Independent Set, Vertex Cover, Dominating Set, colouring, circuit lower bounds) are all
\emph{local} problems.

\subsection{Results}
Throughout, $G=(V,E)$ is a finite simple undirected graph with $n=|V|$ and $m=|E|$, and
$\HC(G)$ is the number of Hamiltonian cycles of $G$, each counted once as a subgraph (so
$\HC(K_4)=3$); $\HC(G)=0$ if $n\le 2$. Running times are counted as in \A: operations on words
long enough to hold an address, plus table queries. All logarithms are to base~$2$.

Our main results are corollaries of \A: they are conditional on three results of \A, which we
recall precisely in \cref{sec:imported}, and they inherit whatever doubt attaches to those
results. If \A{} is wrong, the theorems below lose their proofs.

\begin{hypA}
The following statements of \A{} hold as stated there: Lemma~8, case~(i) (the reduction from
general networks to cliques, \cref{imp:lemma8}); Theorem~A(a),(b) (one algorithm for the
$(K_\ell,\lambda)$-cliques of all levels $\ell$, \cref{imp:thmA}); and Theorem~B (pieces with
parameters $(1/\sqrt2,c)$, \cref{imp:thmB}).
\end{hypA}

For \cref{thm:main-delta} it suffices to have the pieces of \A[Lemma~2] (parameter $11/12$) in
place of Theorem~B, or alternatively the main-text pair \A[Lemma~1, Proposition~3]
(\cref{rem:alt-route}). Nothing else of \A{} is used; in particular we do not use its CNF Main
Theorem, which does not suffice for our purposes (the fermionic sign below is not local at the
level of vertices).

\begin{theorem}[Exact counting, fixed density]\label{thm:main-delta}
Assume $(\mathrm{H}_{\mathrm{A}})$. For every real $\delta>0$ and every $D\ge1$ there is a
deterministic algorithm that, given a graph $G$ with $n$ vertices and at most $Dn$ edges,
outputs $\HC(G)$ within $O(2^{\delta n})$ operations.
\end{theorem}

\begin{theorem}[Exact counting, explicit rate]\label{thm:main-rate}
Assume $(\mathrm{H}_{\mathrm{A}})$. There are an absolute constant $C^{*}$ and one deterministic
algorithm that, given any graph $G$ with $n\ge2$ vertices and $m$ edges, outputs $\HC(G)$
within $2^{C^{*}\max(m,n)/\sqrt{\log n}}$ operations. In particular it runs in time
$2^{O(Dn/\sqrt{\log n})}$ on graphs with $m\le Dn$, where the $O$-constant is absolute and $D$
may depend on $n$.
\end{theorem}

When the density grows with $n$, a variant of the network gives a much better dependence on the
degree. Write $\bar d:=2m/n$ for the average degree and $\Delta$ for the maximum degree.

\begin{theorem}[Exact counting, maximum/average degree; machine-checked]\label{thm:main-degree}
Assume $(\mathrm{H}_{\mathrm{A}})$. There are an absolute constant $C^{\sharp}$ and one
deterministic algorithm that, given any graph $G$ with $n\ge2$ vertices, average degree $\bar d$
and maximum degree $\Delta$, outputs $\HC(G)$ within
\[
  2^{C^{\sharp}n\log(2e\max(\bar d,2))/\sqrt{\log n}}\ \le\ 2^{C^{\sharp}n\log(2e\max(\Delta,2))/\sqrt{\log n}}
\]
operations. In particular it runs in time $2^{o(n)}$ whenever $\log\bar d=o(\sqrt{\log n})$, for
example for $\bar d=(\log n)^{O(1)}$ or $\bar d=2^{(\log n)^{0.49}}$.
\end{theorem}

\Cref{thm:main-degree} is machine-checked: the Lean theorem \lean{theorem3} gives one Turing
machine that outputs $\HC(G)$ in binary within
$\lfloor C\cdot2^{C\,n\log_2(2m/n+2)/\sqrt{\log_2(n+2)}}\rfloor$ steps for every graph $G$ with $n$
vertices and $m$ edges, and \lean{theorem3\_maxdeg} gives the same with $\log_2(\Delta+2)$; both
use the standard axioms only (\cref{sec:verification}). For bounded $\bar d$ it gives the same order
$2^{O(n/\sqrt{\log n})}$ as \cref{thm:main-rate}; it is stronger when $\bar d\to\infty$, because a
pair of blobs then shares only the \emph{set} of used crossing edges between them,
rank-encoded, and the labels of those edges, so the total number of shared bits is about
$2n\log(2e\bar d)$ instead of $2\log_25\cdot m\approx2.32\,\bar dn$. It does not reach dense
graphs: for $\bar d=\Theta(n)$ the bound is $2^{O(n\sqrt{\log n})}$, above $2^n$
(\cref{sec:limitations}).

\begin{corollary}[Decision, search and paths; machine-checked]\label{cor:decision-search}
Assume $(\mathrm{H}_{\mathrm{A}})$. Within time bounds of the same form as
\cref{thm:main-delta,thm:main-rate} (with a larger constant $C$), one can deterministically decide
whether $G$ is Hamiltonian, output a Hamiltonian cycle if one exists, and, for $n\ge2$, count,
decide and find Hamiltonian paths. The search machine makes $n^2+1$ calls to the counting
machine, deleting candidate edges in a fixed order; Hamiltonian paths of $G$ correspond
bijectively to Hamiltonian cycles of $G$ plus an apex vertex.
\end{corollary}
In Lean these are \lean{decision\_TM}, \lean{search\_TM} (fixed density) and
\lean{decision\_TM\_rate}, \lean{search\_TM\_rate} (rate form), and for paths
\lean{hamPaths\_count\_TM}, \lean{hamPaths\_decision\_TM}, \lean{hamPaths\_search\_TM} and their
\lean{\_rate} versions, with the bijection \lean{card\_hamPaths\_eq}; all use the standard axioms
only.

The weighted version is the sparse travelling-salesman problem; it is machine-checked as well.

\begin{theorem}[Sparse TSP; machine-checked]\label{thm:tsp}
Assume $(\mathrm{H}_{\mathrm{A}})$. For every real $\delta>0$ and every $D\ge1$ there are a
deterministic algorithm and a constant $C$ such that, given a graph $G$ with $n$ vertices and at
most $Dn$ edges and integer edge weights $w\colon E\to\{0,\dots,W\}$, the algorithm outputs within
$C\cdot 2^{\delta n}(W+2)^C$ steps the integer $n\sum_{H}2^{b\,w(H)}$, $b=n^2+n$, the sum running
over the Hamiltonian cycles $H$ of $G$. Its base-$2^{b}$ digit $j$ is $n$ times the number of
tours of weight exactly $j$; the minimum tour weight is the position of the lowest non-zero
digit, and the number of optimal tours is that digit divided by $n$. In particular the sparse TSP
is solved exactly, with its optimum counted, in time $2^{\delta n}\poly(W)$ for every
$\delta>0$: subexponential for all weights $W=2^{o(n)}$.
\end{theorem}

In Lean this is \lean{theoremTSP} and \lean{theoremTSP\_decode}, built on \A's development with
the standard axioms only (\cref{sec:verification}). What it does \emph{not} give is as
important: the running time grows with the number of edges, so dense and metric instances
(complete graphs, Euclidean TSP) are untouched, and the $O^*(2^n)$ dynamic programme of
Bellman and of Held and Karp~\cite{Bellman62,HeldKarp62} remains essentially the record there
(improved only to $2^{n-\Omega(\sqrt{n/\log(nW)})}\poly(W)$ for integer weights up to
$W$~\cite{Bjorklund16}, and, for bipartite instances, to $O(1.9999^n)$ assuming quadratic-time
matrix multiplication~\cite{Nederlof20});
and weights $W=2^{\Theta(n)}$ remain open (\cref{sec:limitations}). The identity behind the
theorem is the Kronecker substitution $z=2^{b}$ in $n\sum_Hz^{w(H)}=\sum_Xz^{w(X)}\kappa(X)$,
the weighted form of \cref{thm:blob}.

The next two results are weaker but depend only on the top-level theorems of the labs, which
are themselves machine-checked; the first is machine-checked end to end on top of them.

\begin{theorem}[Parity; machine-checked]\label{thm:parity}
There is one deterministic algorithm that, for every graph $G$ with $n\ge2$ vertices and $m$
edges, computes $\HC(G)\bmod 2$ within $2^{O((m+n)/\sqrt{\log(m+n)})}$ steps, assuming the Main
Theorem of \A. Assuming instead the Lean theorem of~\cite{OpenAI26}, it does so in time
$2^{O(\rho(5m+3n))}$ and polynomial space. Both statements are proved in Lean, and the
hypotheses are discharged by the labs' own Lean proofs, rebuilt from source; the resulting
theorems depend only on the standard axioms \lean{propext}, \lean{Classical.choice},
\lean{Quot.sound}. The analogous statement holds for directed Hamiltonian cycles in digraphs of
bounded in- and out-degree.
\end{theorem}

\begin{theorem}[Randomized decision]\label{thm:randomized}
\leavevmode
\begin{enumerate}[label=\textup{(\alph*)}]
\item \textup{(Machine-checked, assuming the Lean theorem of~\cite{OpenAI26}.)} For every
$\Delta$ and every $\delta>0$ there is a randomized algorithm with one-sided error that decides
Hamiltonicity of graphs of maximum degree $\Delta$ in time $2^{\delta n}$. The same holds for
digraphs of bounded in- and out-degree.
\end{enumerate}
An earlier version also stated a part (b), a randomized $2^{O_D(n(\log\log n)^2/\sqrt{\log n})}$
decision procedure for graphs with $|E|\le Dn$. It is superseded by \cref{thm:main-rate}
(deterministic, and faster) and is no longer claimed; see \cref{rem:randomized-b}.
\end{theorem}

\Cref{thm:parity} and \cref{thm:randomized}(a) are from an earlier stage of this work and are
superseded, as statements, by \cref{thm:main-delta,cor:decision-search}; we include them because
they rest on far less. In (a), Lean checks soundness, the success fraction and the machine; only
the outer loop of independent trials is not modelled.

\Cref{sec:extensions} records further consequences of the same template, for example directed
Hamiltonian cycles, TSP with polynomial weights, Steiner trees, feedback vertex sets and the
Tutte polynomial on the line $y=1$. These are sketched and checked numerically on small
instances; they are \emph{not} fully proved here.

\subsection{Prior work}
No subexponential algorithm for Hamiltonicity on general sparse graphs was known, consistent with
ETH. Planar, bounded-genus and minor-free graphs admit $2^{O(\sqrt n)}$ algorithms via treewidth.
\Cref{tab:sota} lists the best previous bounds in the settings relevant here; all are of the form
$c^n$ with $c>1$. The degree-dependent ones are $(2-\varepsilon)^n$ with $\varepsilon\to0$ as the
degree grows; Björklund's $O^*(1.657^n)$ Monte Carlo decision bound~\cite{Bjorklund14} does not
depend on the degree. The bound of~\cite{BHKK12} is for TSP (hence decision), by a trimmed
Bellman--Held--Karp dynamic programme~\cite{Bellman62,HeldKarp62}; it is not stated for
counting. For small fixed degree, unconditional counting is much faster than
$(2-\varepsilon)^n$: for example, Hamiltonian cycles of graphs of maximum degree~$3$ can be
counted in $O(2^{3n/8})$ time~\cite{Eppstein07}, or in $6^{n/6+o(n)}$ time through the pathwidth
bound $\mathrm{pw}\le n/6+o(n)$ for cubic graphs~\cite{FominHoie06} and the $6^{\mathrm{pw}}$
algorithm of~\cite{BCKN15}.

\begin{table}[ht]
\centering
\small
\caption{Previous bounds and the bounds of this paper (all of the latter conditional; see the
text). Here $\rho(N)=N(\log\log N)^2/\sqrt{\log N}$.}\label{tab:sota}
\begin{tabular}{@{}>{\raggedright\arraybackslash}p{4.2cm}>{\raggedright\arraybackslash}p{5.6cm}>{\raggedright\arraybackslash}p{5.0cm}@{}}
\toprule
Setting & Previously best known & This paper \\
\midrule
undirected, general, decision & $O^{*}(1.657^n)$ randomized~\cite{Bjorklund14}; deterministic $O^{*}(2^n)$~\cite{Bellman62,HeldKarp62,KGK77,Karp82,Bax93}, $2^{n-\Omega(\sqrt{n/\log n})}$~\cite{Bjorklund16} (both also count) & no gain on dense graphs \\
bounded degree $\Delta$, decision/TSP & $(2-\varepsilon_\Delta)^n$~\cite{BHKK12} & deterministic $2^{\delta n}$, exact count \\
maximum degree $3$ & decision/TSP $O(2^{n/3})$~\cite{Eppstein07}, $O^{*}(1.2312^n)$~\cite{XiaoNagamochi16}; counting $O(2^{3n/8})$~\cite{Eppstein07} & deterministic $2^{\delta n}$, exact count \\
counting, pathwidth $\mathrm{pw}$ & $6^{\mathrm{pw}}n^{O(1)}$~\cite[Table~1]{BCKN15}, tight under SETH~\cite[Thm.~1.1]{CLN18} & $2^{O(m/\sqrt{\log n})}$ \\
counting, average degree $\bar d$ & $2^{n-\Omega(n/\bar d)}$ deterministic~\cite{BjorklundWilliams19} (stated for digraphs; covers graphs) & $2^{O(n\log\bar d/\sqrt{\log n})}$ (\cref{thm:main-degree}; machine-checked) \\
counting, directed, average degree $d$ & $2^{n-\Omega(n/d)}$~\cite{BjorklundWilliams19} & see \cref{sec:extensions} (sketch) \\
parity, directed, general & $O^{*}(1.619^n)$~\cite{BjorklundHusfeldt13} & $2^{O(\rho(n))}$ for bounded degree (Lean) \\
parity, undirected, all graphs & -- & $2^{O((m+n)/\sqrt{\log(m+n)})}$ (Lean) \\
\bottomrule
\end{tabular}
\end{table}

Two lines of earlier work are closely related to our technique. Cut\&Count~\cite{CutAndCount11}
reduces connectivity to counting modulo~$2$: a $2$-factor $X$ is paired with the $2^{\cc(X)-1}$
cuts that no edge of $X$ crosses, so non-Hamiltonian $2$-factors cancel in pairs. This is the
route to \cref{thm:parity,thm:randomized}. Björklund's determinant sums~\cite{Bjorklund14} pair
cycle covers with determinants to enforce connectivity; the identity of \cref{lem:2factor-trees}
below is of the same folklore kind. The fermionic, or Grassmann, form of the matrix-tree theorem
that we use is a special case of the generalisation of Kirchhoff's theorem to non-Gaussian
Grassmann integrals by Caracciolo, Jacobsen, Saleur, Sokal and Sportiello~\cite{CJSSS04}. Rank
bounds for the matchings-connectivity matrix (over $\mathbb F_2$~\cite{CKN18}, and over $\Q$~\cite{CLN18},
building on Raz and Spieker's bound for its bipartite variant~\cite{RazSpieker95}) show that, across one cut, connectivity
costs only $O(b)$ bits over $\Q$ rather than $b\log b$; our construction can be read as making
this compose over many cuts simultaneously.

On the hardness side, Dell, Husfeldt, Marx, Taslaman and Wahl\'en~\cite{DHMTW14} and
Curticapean~\cite{Curticapean18} proved $\#$ETH lower bounds for the permanent, perfect matchings
and the Tutte polynomial on sparse graphs, and Brand, Dell and Roth~\cite{BrandDellRoth19}
completed the $\#$ETH dichotomy for the Tutte plane by treating the line $y=1$. Most of these
bounds concern local problems and are already contradicted by the labs' theorems through local
encodings; the line $y=1$ (counting forests) is not local, and we return to it in
\cref{sec:extensions}.

\subsection{Overview of the technique}
\Cref{sec:trees} starts from the identity
\begin{equation}\label{eq:intro-identity}
  n\cdot\HC(G) = \sum_{X\text{ a $2$-factor of } G} \kappa(V,X),
\end{equation}
where $\kappa(V,X)$ is the number of spanning trees of $(V,X)$: a $2$-factor has $n$ spanning
trees if it is a Hamiltonian cycle and none otherwise. The number of spanning trees is the top
coefficient of a product of even elements $1+q_e$ of a Grassmann algebra, one for each edge
(\cref{lem:fermionic-mt}). This replaces the colouring filters of Cut\&Count, which provably
need about $\log n$ bits per vertex to be exact (\cref{prop:colouring-barrier}).

We partition $V$ into $g$ blobs. Each internal edge is handled inside its blob. Each crossing
edge is either unused or used with one of five monomials of $1+q_e$; merging two monomials that
live in the same blob gives five labels, i.e.\ $\log_2 5$ bits per crossing edge
(\cref{tab:labels}). Moving the pieces of the crossing edges into blob order produces a sign. We
show that this sign factors as a product of blob-local signs and an interleaving sign
$(-1)^{\sum_{a<c<b<d}\pi_{ab}\pi_{cd}}$, where $\pi_{ab}$ is the parity of the odd labels on the
bundle of edges between blobs $a$ and $b$ (\cref{prop:sign-rule,prop:sign-factor}). The
interleaving sign is paid for by $O(g^3)$ parity-copy bits (\cref{lem:copies}). The outcome is a
network in the sense of \A{} with arbitrary integer node tables and about $2\log_2 5\cdot m+g^3$
shared bits, whose value is $n\cdot\HC(G)$ (\cref{thm:network-value}). \A's network theorem then
evaluates it (\cref{sec:algorithm}). One trap is that \A's cost depends on the largest number of
bits at one node, so a vertex of degree $\Theta(n)$ must be removed first; we branch on its two
cycle edges (\cref{lem:degree-reduction}). \Cref{sec:degree} replaces the per-edge digits by a
rank code of the set of used crossing edges between two blobs, which costs only logarithmically in
the degree and needs no branching (\cref{thm:main-degree}).

\subsection{What is new here}
Little of the above is new in isolation, and we want to be precise about the credits. The
identity \eqref{eq:intro-identity}, matrix-tree on $2$-factors, is folklore. The fermionic
matrix-tree theorem, which writes the number of spanning trees as a Berezin integral of
$\prod_e(1+q_e)$, is a special case of the generalised matrix-tree theorem of Caracciolo,
Jacobsen, Saleur, Sokal and Sportiello~\cite{CJSSS04}.
Partitioning the vertices into blobs and paying for connectivity across the cuts is the pattern
of Cut\&Count~\cite{CutAndCount11}. Björklund's determinant sums~\cite{Bjorklund14} (FOCS 2010,
SIAM J.\ Comput.\ 2014) pair cycle covers with determinants to enforce connectivity, a closely
related pattern, though there the determinant is evaluated over a field of characteristic~$2$
and the aim is a randomized $O^*(1.657^n)$ bound. The engine that evaluates the final object is
\A's network theorem; we add nothing to it. The one new step is the observation that
connectivity can be encoded with $O(1)$ bits per edge ($\log_25$ per crossing edge) by a
Grassmann determinant whose only non-local part is the fermionic reordering sign, and that this
sign factors so that the non-local remainder costs $O(g^3)$ bits for $g$ blobs. That is exactly
what is needed for the construction to fit the network model of \A, with its linear bit budget
and arbitrary node tables; the Cut\&Count filters do not (\cref{prop:colouring-barrier}), and
neither do the $b\log b$-bit dynamic-programming states over a cut. Everything else is
bookkeeping.

\subsection{Status of the results}
The proofs of \cref{sec:grassmann,sec:trees,sec:blobs,sec:network,sec:algorithm,sec:corollaries}
are complete, conditional on $(\mathrm{H}_{\mathrm{A}})$. They were checked numerically in exact
integer arithmetic, and an independent referee re-derived the construction with fresh code and
found it sound, conditional on \A. The mathematical core is machine-checked in Lean. This
covers the identity $\sum_\ell\prod_a\sigma_a\chi_aI_a=n\cdot\#\mathrm{HC}(G)$ behind
\cref{thm:blob,thm:network-value} and the parameter bounds of \cref{prop:parameters}, with the
network stated in \A's own Lean format. The blob-table procedure is formalised as well, and
\cref{thm:main-delta,thm:main-rate} are both formally verified in Lean end to end, as statements
about Turing machines built on \A's Lean development, with only the standard axioms
(\lean{theorem1}, \lean{theorem2}); a fresh kernel replay with \lean{leanchecker} has passed
for both \lean{theorem1} and \lean{theorem2}. The sparse TSP theorem (\cref{thm:tsp}) is formally
verified at the same level (\lean{theoremTSP}, \lean{theoremTSP\_decode}).
\Cref{thm:main-degree} is the exception: its proof (\cref{sec:degree}) is written out and was
checked numerically and by an independent AI referee, who found it sound, but it is not
formalised. \Cref{sec:verification} gives the details. The bounds are purely asymptotic: as for
the SAT algorithms themselves, the crossover with brute force is astronomically large.

\subsection{Limitations}\label{sec:limitations}
The following should be read together with \cref{sec:verification}.
\begin{enumerate}[label=\textup{(\alph*)}]
\item \emph{Everything is conditional.} The main theorems rest on the claimed ETH refutation
of \A, through $(\mathrm{H}_{\mathrm{A}})$, and on \A's Lean formalisation faithfully capturing
its printed theorems. Expert review of both is pending. Lean guarantees that \A's proofs
establish \A's formal statements, not that those statements say what the printed ones say.
\item \emph{The Lean proofs here were written and reviewed by AI.} They were produced by AI
agents and reviewed by an AI coordinator. No human has yet checked that the formal statements
say what the printed theorems say. A human reader who wants to do so needs to read about five
definitions: \lean{hamCycles} (the set of Hamiltonian cycles being counted, and that it counts
each cycle once as a subgraph); \lean{GraphProblems.code} (the input encoding of a graph, and that
it is the adjacency matrix); \lean{Outputs} (what it means for the machine to output
$\#\mathrm{HC}(G)$ in binary); the statements of \lean{theorem1} and \lean{theorem2} themselves,
in particular that the step bound is $\lfloor C\cdot2^{\delta n}\rfloor$ resp.\
$\lfloor C\cdot2^{C(m+n)/\sqrt{\log_2(n+2)}}\rfloor$ with $C$ independent of $G$; and \A's
\lean{Net} format together with the hypothesis bundle through which our network is handed to
\A's theorems, to confirm that nothing beyond \cref{imp:lemma8,imp:thmA,imp:thmB} is assumed.
\item \emph{Sparse graphs only.} The network of \cref{sec:network} has about $2\log_25$ bits per
edge, so \cref{thm:main-rate} beats $2^n$ only when $m=o(n\sqrt{\log n})$. The rank-encoded
network of \cref{sec:degree} has about $2\log(2e\bar d)$ bits per vertex, so
\cref{thm:main-degree} extends this, roughly, to average degree $\bar d=2^{o(\sqrt{\log n})}$. Dense
graphs are still not covered: at $\bar d\approx\Delta\approx n$ the bound of
\cref{thm:main-degree} is $2^{\Theta(n\log n/\sqrt{\log n})}$, which exceeds $2^n$. General TSP
and the Held--Karp $O^*(2^n)$ bound are untouched.
\item \emph{``Previous algorithms were exponential'' refers to general sparse graphs.} Planar,
bounded-genus, minor-free and bounded-treewidth graphs already had subexponential or
single-exponential-in-treewidth algorithms, and on planar graphs the $2^{O(\sqrt n)}$ bound is
better than ours. The gain is on sparse graphs with no such structure, for example bounded-degree
expanders.
\item \emph{Asymptotic only.} The crossover with brute force is astronomically large; nothing
here is of practical use.
\item \emph{The extensions are sketched.} \Cref{sec:extensions} is checked numerically on
small instances and is not proved at the level of the rest of the paper.
\item \emph{Citations.} All citations were checked against the primary sources on 30 Oct 2026.
In particular, the $6^{\mathrm{pw}}n^{O(1)}$ counting bound in \cref{tab:sota} is
$\tilde O(6^{\mathrm{pw}}\mathrm{pw}^{O(1)}n^2)$ in \cite[Table~1]{BCKN15}, and its SETH-tightness
is \cite[Thm.~1.1]{CLN18}; the other constants in that table are as stated in the cited works.
\item \emph{No systematic novelty search has been done.} The credits in the previous subsection
are to the works we know; a Grassmann-determinant encoding of connectivity in a bit-bounded model
may well exist elsewhere, and we have not searched the literature systematically.
\end{enumerate}

\subsection{Organisation}
\Cref{sec:imported} recalls the results of \A{} that we use. \Cref{sec:grassmann} collects facts
about Grassmann algebras. \Cref{sec:trees} proves the fermionic matrix-tree identity for
Hamiltonian cycles. \Cref{sec:blobs} splits it over blobs and proves the sign rule.
\Cref{sec:network} builds the network and bounds its parameters. \Cref{sec:algorithm} proves
\cref{thm:main-delta,thm:main-rate}, \cref{sec:degree} proves \cref{thm:main-degree} with a
rank-encoded variant of the network, and \cref{sec:corollaries} proves
\cref{cor:decision-search}. \Cref{sec:parity} treats the results that rest only on the top-level
theorems. \Cref{sec:extensions} sketches extensions, \cref{sec:verification} describes the
verification status, and \cref{sec:open} lists open problems.

% =====================================================================
\section{The network theorem of \texorpdfstring{\A}{[A]}}\label{sec:imported}
% =====================================================================

We recall the definitions and results of \A{} that we use, in the notation of \A{} except where
a letter clashes with ours. Conventions of \A: $\log=\log_2$, $[x=u]\in\{0,1\}$ and
$[K]=\{1,\dots,K\}$.

\begin{definition}[{\A[Definition~1]}]\label{def:factor}
A function $F\colon X\times X'\to\Z$ \emph{factors through $c$ bits} if there are
$f\colon X\times\{0,1\}^c\to\Z$ and $f'\colon X'\times\{0,1\}^c\to\Z$ with
$F(z,z')=\sum_{\theta\in\{0,1\}^c} f(z,\theta)f'(z',\theta)$. In particular, for any
$h\colon X\to\{0,1\}^c$ and $h'\colon X'\to\{0,1\}^c$ the function
$[h(z)=h'(z')]=\sum_\theta [h(z)=\theta][h'(z')=\theta]$ factors through $c$ bits.
\end{definition}

\begin{definition}[{\A[Definition~2]}]\label{def:network}
Let $P,B,\tau\ge1$ be integers. A \emph{network on $J$ with parameters $P,B,\tau$} consists of a
finite set $J$ of nodes; for each node $a$ a finite set $S_a$ of settings and a function
$W_a\colon S_a\to\Z$; and for each unordered pair $\{a,b\}$ of distinct nodes a function
$\Psi_{a,b}\colon S_a\times S_b\to\Z$ (with $\Psi_{b,a}(z_b,z_a)=\Psi_{a,b}(z_a,z_b)$), an integer
$c_{a,b}=c_{b,a}\ge0$ and a factorisation $(f_{a,b},f_{b,a})$ of $\Psi_{a,b}$ through $c_{a,b}$
bits. The settings of $a$ are encoded as bit strings of one fixed length $P_a\le P$, and
\begin{enumerate}[label=\textup{(\roman*)}]
\item $\sum_{b\ne a} c_{a,b}\le P$ for every $a\in J$;
\item all values of all $W_a$, $\Psi_{a,b}$, $f_{a,b}$ are below $2^B$ in absolute value;
\item an algorithm receives $P$, $B$ and the numbers $P_a$, $c_{a,b}$ as input, and can decide
membership in $S_a$ and compute any value of $W_a$, $\Psi_{a,b}$, $f_{a,b}$ with at most $\tau$
operations.
\end{enumerate}
The \emph{value} of the network is
$Z=\sum_{z\in\prod_a S_a}\prod_{a\in J}W_a(z_a)\prod_{\{a,b\}}\Psi_{a,b}(z_a,z_b)$, and its
\emph{total} is $U=\sum_{(a,b)\colon a\ne b}c_{a,b}$, which counts every shared bit from both
ends. By (i), $U\le |J|P$.
\end{definition}

\begin{definition}[{\A[Definition~3]}]
Let $K\ge2$ and $\lambda,B\ge1$. In the \emph{$(K,\lambda)$-clique} with tables
$W_1,\dots,W_K$, a setting of node $w$ is a tuple $(z_w(w'))_{w'\ne w}$ of $\lambda$-bit strings,
$c_{w,w'}=\lambda$ and $\Psi_{w,w'}=[z_w(w')=z_{w'}(w)]$. The tables $W_w$, with $|W_w|<2^B$, are
queried at unit cost.
\end{definition}

\begin{importedlemma}[{\A[Lemma~8, case~(i)]}]\label{imp:lemma8}
Let $K\ge2$, $g=3K^8$ and $\zeta=1/(2K)$. Take a network on $[g]$ with parameters $P,B,\tau$ and
total $U$, and put
\[
  L=\Bigl\lceil\frac{U+gP/K}{K^2}\Bigr\rceil,\qquad B'=\lceil 2gP/K\rceil+g^2B .
\]
If an algorithm outputs the value of every $(K,L)$-clique with $|W_w|<2^{B'}$ within $T$
operations, then the value of the network is computed within
\begin{equation}\label{eq:lemma8}
  \bigl(K^g+T\cdot 2^{2gP/K}\cdot\tau\bigr)\cdot\bigl(g(P+B)\bigr)^{C_1}
\end{equation}
operations, by one procedure, for an absolute constant $C_1$.
\end{importedlemma}

The bound~\eqref{eq:lemma8} is weaker than the printed form of \A[Lemma~8] by at most a
polynomial factor, which is irrelevant below; it is the form proved in the Lean development
of \A{} as \lean{ExplicitCost.groupingExplicit}.

\begin{importedtheorem}[{\A[Theorem~A]}]\label{imp:thmA}
Suppose there are pieces with parameters $(\rho,c)$ in the sense of \A[Definition~7], and let
$\alpha'=\log(1/\rho)$. Fix an integer $x_0\ge18$ with $2^{x_0}\ge c+5$ and define
\[
  K_0=2^{x_0},\quad \zeta_\ell=2^{-\ell-5},\quad g_\ell=2^{3\ell+x_0}K_\ell,\quad K_{\ell+1}=g_\ell^2,\quad
  e_0=\tfrac12,\quad e_{\ell+1}=(\rho+3\zeta_\ell)e_\ell+3/K_\ell .
\]
\begin{enumerate}[label=\textup{(\alph*)}]
\item $\log K_\ell=(3x_0+6)2^\ell-6\ell-2x_0-6$, so $K_\ell<K_{\ell+1}\le K_\ell^{10}$; and
$\frac12\rho^\ell\le e_\ell\le 0.85\,\rho^\ell$, where $\rho^\ell\le((3x_0+6)/\log K_\ell)^{\alpha'}$.
\item One deterministic algorithm, given $\ell\ge0$, $\lambda,B\ge1$ and query access to tables
with $|W_w|<2^B$, outputs the value of the $(K_\ell,\lambda)$-clique within
\begin{equation}\label{eq:thmA}
  2^{e_\ell K_\ell^2\lambda}\cdot\bigl(K_\ell(\lambda+B)\bigr)^{C_0K_\ell^2}
\end{equation}
operations, with $C_0$ depending only on $c$.
\end{enumerate}
\end{importedtheorem}

\begin{importedtheorem}[{\A[Theorem~B and Lemma~2]}]\label{imp:thmB}
There are pieces with parameters $(1/\sqrt2,c)$ for an absolute constant $c$. Moreover,
\A[Lemma~2] gives pieces with parameters $(11/12,c)$ for every integer $c\ge 8$ large enough.
\end{importedtheorem}

\begin{remark}[Alternative route for \cref{thm:main-delta}]\label{rem:alt-route}
\A[Lemma~1] states: if $(K,e)$ is achievable and $g=3K^8$, every network on $[g]$ is evaluated
within $C_1 2^{e(U+gP/K)+2gP/K}(P+B)^{C_1}\tau$ operations, with $C_1=C_1(K,e)$; and
\A[Proposition~3] gives, for every $\eta>0$, infinitely many $K$ with $(K,\eta)$ achievable. This
pair can replace $(\mathrm{H}_{\mathrm{A}})$ in the proof of \cref{thm:main-delta}.
\end{remark}

\begin{remark}
\A's Lean development proves \cref{imp:lemma8} and \cref{imp:thmA}(b) as
\lean{ExplicitCost.groupingExplicit} and \lean{ExplicitCost.recursionSolves}, and \cref{imp:thmB}
as \lean{PiecesCounted.halfPieces}. There they are instantiated only for the networks that come
from CNFs, and Lean's machine model differs from the printed one (operations are polynomial-time
steps or queries, memory is bounded, and $K,g,P,B$ are given in unary). Applying these results to
the network of \cref{sec:network} is done here on paper only; see \cref{sec:verification}.
\end{remark}

% =====================================================================
\section{Grassmann algebra}\label{sec:grassmann}
% =====================================================================

Let $\Gamma$ be a finite set of \emph{generators} with a total order $\prec$. The exterior algebra
$\Lambda(\Gamma)$ over $\Z$ is the free $\Z$-module with basis $\{e_S : S\subseteq\Gamma\}$ and
product
\[
  e_S\cdot e_T=0\ \text{ if } S\cap T\ne\emptyset,\qquad
  e_S\cdot e_T=(-1)^{\#\{(s,t)\in S\times T\,:\,t\prec s\}}\,e_{S\cup T}\ \text{ otherwise}.
\]
This is the exterior algebra of $\Z^\Gamma$~\cite[Ch.~III, \S7]{Bourbaki89}; it is associative
with unit $e_\emptyset=1$. Writing $\gamma$ for $e_{\{\gamma\}}$ we have $\gamma\gamma'=-\gamma'\gamma$
and $\gamma^2=0$, and $e_S$ is the product of the elements of $S$ in increasing $\prec$-order. A
product of generators therefore does not depend on $\prec$, and all statements below, being about
sums and products of generators, are independent of $\prec$; changing $\prec$ only re-signs the
basis.

\emph{Parity.} $\Lambda=\Lambda^0\oplus\Lambda^1$, where $\Lambda^p$ is spanned by the $e_S$ with
$|S|\equiv p\pmod 2$. Elements of $\Lambda^0\cup\Lambda^1$ are \emph{homogeneous} (so $0$ is
homogeneous of both parities), and $\Lambda^0$ is a subring.

\begin{fact}[Graded commutativity]\label{fact:graded}
For $x\in\Lambda^p$ and $y\in\Lambda^q$ we have $xy=(-1)^{pq}yx$. In particular $\Lambda^0$ is
central.
\end{fact}

\begin{proof}
For disjoint $S,T$, every pair in $S\times T$ is counted in exactly one of $\#\{t\prec s\}$ and
$\#\{s\prec t\}$, so $e_Se_T=(-1)^{|S||T|}e_Te_S$. Extend bilinearly.
\end{proof}

\begin{lemma}[Koszul sign rule]\label{lem:koszul}
Let $y_1,\dots,y_M$ be homogeneous of parities $p_1,\dots,p_M$ and let $\sigma$ be a permutation
of $[M]$. Then
\[
  y_{\sigma(1)}y_{\sigma(2)}\cdots y_{\sigma(M)}=(-1)^{N(\sigma)}\,y_1y_2\cdots y_M,\qquad
  N(\sigma)=\sum_{\substack{s<t\\ \sigma^{-1}(s)>\sigma^{-1}(t)}}p_sp_t,
\]
the sum running over the pairs whose relative order $\sigma$ reverses.
\end{lemma}

\begin{proof}
Induction on the number of reversed pairs. If $\sigma\ne\mathrm{id}$, some adjacent positions
$j,j+1$ have $\sigma(j)>\sigma(j+1)$. Exchanging these two factors multiplies the product by
$(-1)^{p_{\sigma(j)}p_{\sigma(j+1)}}$ by \cref{fact:graded}; it removes exactly the pair
$\{\sigma(j),\sigma(j+1)\}$ from the set of reversed pairs and changes no other relative order.
\end{proof}

\emph{Subalgebras and top coefficients.} For $\Gamma'\subseteq\Gamma$, $\Lambda(\Gamma')\subseteq
\Lambda(\Gamma)$ is the span of the $e_S$ with $S\subseteq\Gamma'$. Let $|\Gamma'|$ be even and
$\omega=\pm e_{\Gamma'}$. For $x\in\Lambda(\Gamma')$ we write $\topc{\omega}{x}$ for the coefficient
$c$ such that $x-c\,\omega$ lies in the span of the $e_S$ with $S\subsetneq\Gamma'$. (This is the
Berezin integral of $x$, normalised by $\omega$.)

\begin{lemma}[Block factorisation of the top coefficient]\label{lem:block}
Let $\Gamma=\Gamma_1\sqcup\dots\sqcup\Gamma_g$ with every $|\Gamma_a|$ even (blocks may be empty).
Let $\omega_a=\pm e_{\Gamma_a}$, $\Omega=\omega_1\omega_2\cdots\omega_g$ and
$x_a\in\Lambda(\Gamma_a)$. Then
\[
  \topc{\Omega}{x_1x_2\cdots x_g}=\prod_{a=1}^{g}\topc{\omega_a}{x_a}.
\]
\end{lemma}

\begin{proof}
Since $\Omega=\pm e_\Gamma$, $\topc{\Omega}{\cdot}$ is defined. Write
$x_a=\topc{\omega_a}{x_a}\,\omega_a+r_a$, where $r_a$ is a combination of $e_S$ with
$S\subsetneq\Gamma_a$. Expanding the product, the term with no $r_a$ is
$\prod_a\topc{\omega_a}{x_a}\cdot\Omega$. Every other term is a combination of products
$e_{S_1}\cdots e_{S_g}$ with $S_b\subseteq\Gamma_b$ for all $b$ and $S_a\subsetneq\Gamma_a$ for
some $a$; such a product is $0$ or $\pm e_{\bigcup S_b}$ with $\bigcup S_b\ne\Gamma$, and does not
contribute to $\topc{\Omega}{\cdot}$.
\end{proof}

% =====================================================================
\section{Spanning trees and Hamiltonian cycles}\label{sec:trees}
% =====================================================================

Fix $G=(V,E)$ with $n\ge3$ and a root $r\in V$. A \emph{$2$-factor} is a set $X\subseteq E$ such
that every vertex has degree exactly $2$ in $(V,X)$. We write $\kappa(V,X)$ for the number of
spanning trees of the graph $(V,X)$.

\begin{lemma}\label{lem:2factor-trees}
$n\cdot\HC(G)=\sum_{X}\kappa(V,X)$, the sum over all $2$-factors $X$ of $G$.
\end{lemma}

\begin{proof}
A $2$-factor $X$ makes $(V,X)$ a disjoint union of cycles covering $V$. If $(V,X)$ is connected
it is a single $n$-cycle, i.e.\ a Hamiltonian cycle, whose spanning trees are obtained by deleting
one of its $n$ edges; so $\kappa(V,X)=n$. If $(V,X)$ is disconnected, $\kappa(V,X)=0$. Conversely,
the edge set of a Hamiltonian cycle is a connected $2$-factor, and distinct Hamiltonian cycles
have distinct edge sets.
\end{proof}

\emph{Generators.} Let $\Gamma=\{\bpsi_v,\psi_v : v\in V\setminus\{r\}\}$, in any order. We put
$\bpsi_r=\psi_r:=0$ (the zero element), $\omega_v:=\bpsi_v\psi_v\in\Lambda^0$, and for $S\subseteq
V$ and $e=\{i,j\}\in E$,
\[
  \Omega_S:=\prod_{v\in S\setminus\{r\}}\omega_v,\qquad
  q_e:=(\bpsi_i-\bpsi_j)(\psi_i-\psi_j)\in\Lambda^0 .
\]
The order of the factors of $\Omega_S$ is irrelevant since they are even and central, and
$\Omega_S=\pm e_{\Gamma_S}$, where $\Gamma_S$ is the set of generators of the vertices of $S$. The
element $q_e$ does not depend on the orientation of $e$, since both factors change sign.

\begin{lemma}[Berezin form of Cauchy--Binet]\label{lem:berezin-cb}
Let $N=n-1$ and $b_1,\dots,b_N\in\Z^{V\setminus\{r\}}$. Put $\bar\theta(b)=\sum_u b(u)\bpsi_u$ and
$\theta(b)=\sum_u b(u)\psi_u$. Then
\[
  \topc{\Omega_V}{\bar\theta(b_1)\theta(b_1)\,\bar\theta(b_2)\theta(b_2)\cdots\bar\theta(b_N)\theta(b_N)}=\det(\mathbf B)^2,
\]
where $\mathbf B$ is the $N\times N$ matrix with columns $b_k$, its rows indexed by
$V\setminus\{r\}$ in any fixed order; $\det(\mathbf B)^2$ does not depend on that order.
\end{lemma}

\begin{proof}
Enumerate $V\setminus\{r\}$ as $u_1,\dots,u_N$.
\begin{enumerate}[label=(\arabic*)]
\item By \cref{lem:koszul}, moving each $\bar\theta(b_k)$ to the left past
$\theta(b_1),\dots,\theta(b_{k-1})$ gives
$\prod_k\bar\theta(b_k)\theta(b_k)=(-1)^{N(N-1)/2}\,\bar\theta(b_1)\cdots\bar\theta(b_N)\cdot
\theta(b_1)\cdots\theta(b_N)$.
\item Expanding multilinearly,
$\bar\theta(b_1)\cdots\bar\theta(b_N)=\sum_{\varphi\colon[N]\to[N]}\prod_k b_k(u_{\varphi(k)})\;
\bpsi_{u_{\varphi(1)}}\cdots\bpsi_{u_{\varphi(N)}}$. Only bijections $\varphi$ survive, because
$\bpsi_u^2=0$, and for those $\bpsi_{u_{\varphi(1)}}\cdots\bpsi_{u_{\varphi(N)}}=\operatorname{sgn}
(\varphi)\,\bpsi_{u_1}\cdots\bpsi_{u_N}$ by \cref{lem:koszul}. Hence
\[
  \bar\theta(b_1)\cdots\bar\theta(b_N)=\det(\mathbf B)\,\bpsi_{u_1}\cdots\bpsi_{u_N},
  \qquad
  \theta(b_1)\cdots\theta(b_N)=\det(\mathbf B)\,\psi_{u_1}\cdots\psi_{u_N}.
\]
\item Step (1) applied to the standard vectors gives
\[
  \bpsi_{u_1}\cdots\bpsi_{u_N}\,\psi_{u_1}\cdots\psi_{u_N}=(-1)^{N(N-1)/2}\,\omega_{u_1}\cdots
  \omega_{u_N}=(-1)^{N(N-1)/2}\,\Omega_V .
\]
\end{enumerate}
The two signs cancel and the product equals $\det(\mathbf B)^2\,\Omega_V$.
\end{proof}

\begin{lemma}[Fermionic matrix-tree theorem~\cite{CJSSS04}]\label{lem:fermionic-mt}
For every $X\subseteq E$,
\[
  \topc{\Omega_V}{\prod_{e\in X}(1+q_e)}=\kappa(V,X).
\]
\end{lemma}

\begin{proof}
The factors are even and commute, so $\prod_{e\in X}(1+q_e)=\sum_{F\subseteq X}\prod_{e\in F}q_e$.
The term of $F$ lies in the span of the $e_S$ with $|S|=2|F|$, while $\Omega_V=\pm e_\Gamma$ with
$|\Gamma|=2N$; so only $|F|=N$ contributes. For $e=\{i,j\}$ we have
$q_e=\bar\theta(b_e)\theta(b_e)$ with $b_e=\mathbf 1_i-\mathbf 1_j$ restricted to
$V\setminus\{r\}$ (that is, $b_e=\mathbf 1_i$ if $j=r$). By \cref{lem:berezin-cb}, $F$ contributes
$\det(\mathbf B_F)^2$, where $\mathbf B_F$ is the reduced incidence matrix of $F$.
\begin{itemize}
\item If $F$ (with $|F|=n-1$) is not a spanning tree, then $(V,F)$ is disconnected, since a
connected graph on $n$ vertices with $n-1$ edges is a tree. Let $C$ be a component with $r\notin
C$. The rows $u\in C$ of $\mathbf B_F$ sum to $0$: the column of an edge inside $C$ has entries
$+1$ and $-1$ there, and the other columns vanish on $C$. So $\det\mathbf B_F=0$.
\item If $F$ is a spanning tree, then $\det\mathbf B_F=\pm1$, by induction on $n$. For $n=1$ the
matrix is empty. For $n\ge2$ pick a leaf $v\ne r$ (a tree has at least two leaves). Row $v$ has
one nonzero entry, $\pm1$, in the column of its edge $f$, and expanding along it gives
$\det\mathbf B_F=\pm\det\mathbf B_{F-f}$, the reduced incidence matrix of the spanning tree
$F-f$ of $V-v$.
\end{itemize}
Hence the coefficient is the number of spanning trees $F\subseteq X$, i.e.\ $\kappa(V,X)$.
\end{proof}

By Cauchy--Binet, \cref{lem:fermionic-mt} is equivalent to Kirchhoff's theorem
$\kappa(V,X)=\det L_X^{(r)}$; we do not need this form, but it is checked numerically
(\cref{sec:verification}).

\begin{corollary}\label{cor:hc-berezin}
\[
  n\cdot\HC(G)=\sum_{x\in\{0,1\}^E}\ \prod_{v\in V}\Ind{\deg_x(v)=2}\cdot
  \topc{\Omega_V}{\prod_{e\colon x_e=1}(1+q_e)} .
\]
\end{corollary}

\begin{proof}
Combine \cref{lem:2factor-trees,lem:fermionic-mt}.
\end{proof}

% =====================================================================
\section{Blobs, pieces and the sign rule}\label{sec:blobs}
% =====================================================================

\subsection{Conventions}
The following conventions are fixed once and for all.
\begin{itemize}
\item \emph{Partition.} $\phi\colon V\to[g]$ is an arbitrary map, and the \emph{blobs} are
$V_a=\phi^{-1}(a)$. Blobs may be empty, and $r$ may lie in any blob. The blob indices
$1<2<\dots<g$ form the ``line'' on which chords interleave.
\item \emph{Edges.} An edge is \emph{internal} to $a$ if both ends lie in $V_a$; the set of these
is $\Eint(a)$. The other edges are \emph{crossing}; let $\mx$ be their number.
\item \emph{Crossing order.} The crossing edges are listed in an arbitrary but fixed total order
$e_1,\dots,e_{\mx}$; we write $k<l$ if $e_k$ comes before $e_l$.
\item \emph{Ends.} A crossing edge $e_k$ has a \emph{low end} $s_k$ and a \emph{high end} $t_k$,
with $a_k:=\phi(s_k)<b_k:=\phi(t_k)$. Its \emph{chord} is $\{a_k,b_k\}$.
\item \emph{Incidences.} $\partial(a):=\{k : a\in\{a_k,b_k\}\}$ is the set of crossing edges at
$a$. For $a<b$ let $E(a,b):=\{k : (a_k,b_k)=(a,b)\}$ (the \emph{bundle} of $\{a,b\}$) and
$m_{ab}:=|E(a,b)|$, and set $m_{ba}:=m_{ab}$.
\item \emph{Algebras.} $\Gamma_a$ is the set of generators of $V_a\setminus\{r\}$,
$\Lambda_a:=\Lambda(\Gamma_a)$ and $\omega_a:=\Omega_{V_a}$. Then $\Omega_V=\omega_1\cdots\omega_g$,
and \cref{lem:block} applies to $\Gamma=\bigsqcup_a\Gamma_a$.
\end{itemize}

\subsection{Labels and pieces}
The label set is $\cL=\{\lab N,\lab 0,\lab 1,\lab 2,\lab 3\}$, and the \emph{parity} of a label is
$\pi(\ell):=\Ind{\ell\in\{\lab 2,\lab 3\}}$. For the crossing edge $e_k$ with ends $s=s_k$ and
$t=t_k$, the pieces $\alpha_\ell(k)\in\Lambda_{a_k}$ and $\beta_\ell(k)\in\Lambda_{b_k}$ are given
in \cref{tab:labels}.

\begin{table}[ht]
\centering
\caption{Labels of a crossing edge $e_k$ with low end $s$ and high end $t$, and their pieces.}
\label{tab:labels}
\begin{tabular}{@{}cccccl@{}}
\toprule
label $\ell$ & $\alpha_\ell(k)\in\Lambda_{a_k}$ & $\beta_\ell(k)\in\Lambda_{b_k}$ & $\alpha_\ell(k)\beta_\ell(k)$ & parity $\pi(\ell)$ & meaning \\
\midrule
$\lab N$ & $1$ & $1$ & $1$ & $0$ & $e_k\notin X$\\
$\lab 0$ & $1$ & $1+\bpsi_t\psi_t$ & $1+\bpsi_t\psi_t$ & $0$ & $e_k\in X$, terms $1$ and $\bpsi_t\psi_t$\\
$\lab 1$ & $\bpsi_s\psi_s$ & $1$ & $\bpsi_s\psi_s$ & $0$ & $e_k\in X$, term $\bpsi_s\psi_s$\\
$\lab 2$ & $\bpsi_s$ & $-\psi_t$ & $-\bpsi_s\psi_t$ & $1$ & $e_k\in X$, term $-\bpsi_s\psi_t$\\
$\lab 3$ & $\psi_s$ & $\bpsi_t$ & $\psi_s\bpsi_t$ & $1$ & $e_k\in X$, term $-\bpsi_t\psi_s$\\
\bottomrule
\end{tabular}
\end{table}

\begin{remark}[Why five labels]\label{rem:five-labels}
For a used edge, $1+q_e=1+\bpsi_s\psi_s-\bpsi_s\psi_t-\bpsi_t\psi_s+\bpsi_t\psi_t$ has five
monomials; together with ``unused'' this gives six states. The two monomials $1$ and
$\bpsi_t\psi_t$ both live entirely in the high blob, and label $\lab 0$ merges them into the
single high piece $1+\bpsi_t\psi_t$. This gives five labels, i.e.\ $\log_25$ bits per crossing
edge. (A description as ``unused, or one of four bilinear terms'' would miss the term $1$; the
merge above is what makes five labels correct. Other merges work equally well and change only
constants.)
\end{remark}

\begin{lemma}\label{lem:pieces}
\leavevmode
\begin{enumerate}[label=\textup{(\alph*)}]
\item $\alpha_\ell(k)$ and $\beta_\ell(k)$ are homogeneous of parity $\pi(\ell)$. This includes
the cases $s=r$ or $t=r$, in which some pieces are $0$ or reduce to $1$.
\item $\alpha_{\lab N}(k)\beta_{\lab N}(k)=1$ and
$\sum_{\ell\in\{\lab0,\lab1,\lab2,\lab3\}}\alpha_\ell(k)\beta_\ell(k)=1+q_{e_k}$.
\end{enumerate}
\end{lemma}

\begin{proof}
(a) Read off \cref{tab:labels}. (b) By bilinearity
$q_e=\bpsi_s\psi_s-\bpsi_s\psi_t-\bpsi_t\psi_s+\bpsi_t\psi_t$ for any elements, zero included, and
$-\bpsi_t\psi_s=\psi_s\bpsi_t$. The four products in the table are $1+\bpsi_t\psi_t$,
$\bpsi_s\psi_s$, $-\bpsi_s\psi_t$ and $\psi_s\bpsi_t$, which sum to $1+q_e$.
\end{proof}

For a labelling $\eta\in\cL^{\mx}$ write $p_k:=\pi(\eta_k)$. Define the \emph{piece of $k$ at
$a$} as $\gamma_{k,a}(\eta):=\alpha_{\eta_k}(k)$ if $a=a_k$ and $\beta_{\eta_k}(k)$ if $a=b_k$,
and the ordered product of the pieces at $a$,
\[
  R_a(\eta):=\prod_{k\in\partial(a)}^{\to}\gamma_{k,a}(\eta)\in\Lambda_a
  \qquad\text{(factors in increasing crossing order $k$).}
\]
Finally, let
\[
  r_{kl}:=\#\bigl\{(x,y)\in\{a_k,b_k\}\times\{a_l,b_l\} : y<x\bigr\},\qquad
  \varepsilon(\eta):=(-1)^{\sum_{k<l}p_kp_l\,r_{kl}} .
\]

\begin{proposition}[Sign rule]\label{prop:sign-rule}
For every partition, every crossing order and every $\eta\in\cL^{\mx}$,
\[
  \prod_{k=1}^{\mx}\alpha_{\eta_k}(k)\,\beta_{\eta_k}(k)=\varepsilon(\eta)\cdot R_1(\eta)R_2(\eta)\cdots R_g(\eta).
\]
\end{proposition}

\begin{proof}
The left side is the product of the sequence
$Y=(\alpha(1),\beta(1),\alpha(2),\beta(2),\dots,\alpha(\mx),\beta(\mx))$. Give the piece
$\alpha(k)$ the key $(a_k,k)$ and the piece $\beta(k)$ the key $(b_k,k)$. Sorting $Y$ by key,
lexicographically, lists blob~$1$'s pieces in increasing $k$, then blob~$2$'s, and so on; so the
sorted product is $R_1\cdots R_g$. By \cref{lem:pieces}(a) all pieces are homogeneous, and
\cref{lem:koszul} counts the pairs of pieces whose relative order is reversed:
\begin{itemize}
\item \emph{Two pieces of the same $k$.} In $Y$, $\alpha(k)$ comes first, and its key is also
smaller since $a_k<b_k$. Not reversed.
\item \emph{A piece $x$ of $k$ and a piece $y$ of $l$, with $k<l$.} In $Y$, $x$ comes first. The
keys are $(\mathrm{blob}(x),k)$ and $(\mathrm{blob}(y),l)$; since $k<l$, the pair is reversed iff
$\mathrm{blob}(y)<\mathrm{blob}(x)$. The product of parities is $p_kp_l$, and the number of such
pairs is $r_{kl}$.
\end{itemize}
So $R_1\cdots R_g=(-1)^{\sum_{k<l}p_kp_lr_{kl}}\prod_k\alpha(k)\beta(k)$; multiply both sides by
$\varepsilon(\eta)=\pm1$.
\end{proof}

\subsection{Factorisation of the sign}

\begin{lemma}\label{lem:disjoint-chords}
If $\{a_k,b_k\}\cap\{a_l,b_l\}=\emptyset$, then $r_{kl}\equiv I(k,l)\pmod 2$, where
\[
  I(k,l):=\Ind{a_k<a_l<b_k<b_l\ \text{or}\ a_l<a_k<b_l<b_k}
\]
indicates that the chords interleave. In particular $r_{kl}\bmod 2$ is symmetric and independent
of the crossing order.
\end{lemma}

\begin{proof}
Put $Y=\{a_l,b_l\}$, whose elements differ from $a_k$ and $b_k$. Then
$r_{kl}=\#\{y\in Y: y<a_k\}+\#\{y\in Y: y<b_k\}=2\#\{y\in Y: y<a_k\}+\#\{y\in Y: a_k<y<b_k\}$. The
last number is odd iff exactly one endpoint of chord $l$ lies strictly between $a_k$ and $b_k$,
that is, iff the chords interleave.
\end{proof}

Two crossing edges $k<l$ \emph{share a blob} if their chords meet. For such a pair put
$s(k,l):=\min(\{a_k,b_k\}\cap\{a_l,b_l\})$; then $k,l\in\partial(s(k,l))$. Define the \emph{local
signs} and the \emph{bundle parities}
\[
  \sigma_a(\eta):=(-1)^{\sum_{k<l\,:\ \text{chords meet},\ s(k,l)=a}p_kp_l\,r_{kl}},\qquad
  \pi_{ab}(\eta):=\sum_{k\in E(a,b)}p_k\bmod 2\quad(a<b).
\]
Thus $\sigma_a$ is a function of $\eta|_{\partial(a)}$, and $\pi_{ab}$ of $\eta|_{E(a,b)}$.

\begin{proposition}[Sign factorisation]\label{prop:sign-factor}
\[
  \varepsilon(\eta)=\prod_{a=1}^{g}\sigma_a(\eta)\cdot(-1)^{\sum_{a<c<b<d}\pi_{ab}(\eta)\,\pi_{cd}(\eta)} .
\]
\end{proposition}

\begin{proof}
Split $\sum_{k<l}p_kp_lr_{kl}$ into the pairs whose chords meet and the pairs whose chords are
disjoint.
\begin{itemize}
\item \emph{Meeting chords.} Each such pair has exactly one $s(k,l)$, so these pairs give
$\prod_a\sigma_a$.
\item \emph{Disjoint chords.} By \cref{lem:disjoint-chords} these pairs give, modulo $2$,
$\sum_{k<l\colon I(k,l)=1}p_kp_l$. Take an unordered pair $\{k,l\}$ with disjoint interleaving
chords, let $\{a,b\}$ ($a<b$) be the chord with the smaller minimum and $\{c,d\}$ the other;
interleaving means $a<c<b<d$. Conversely, every quadruple $a<c<b<d$ with $k\in E(a,b)$ and $l\in
E(c,d)$ gives such a pair. This is a bijection, so the sum equals
$\sum_{a<c<b<d}\bigl(\sum_{k\in E(a,b)}p_k\bigr)\bigl(\sum_{l\in E(c,d)}p_l\bigr)\equiv
\sum_{a<c<b<d}\pi_{ab}\pi_{cd}\pmod2$.\qedhere
\end{itemize}
\end{proof}

\emph{Parity copies.} Let $Q_a:=\{(c,d) : a<c<d\le g\}$, so $|Q_a|=\binom{g-a}{2}$. For
$y\in\{0,1\}^{Q_a}$ put
\[
  \chi_a(\eta,y):=(-1)^{\sum_{(c,d)\in Q_a}y_{cd}\sum_{b\colon c<b<d}\pi_{ab}(\eta)},
\]
a function of $\eta|_{\partial(a)}$ and $y$.

\begin{lemma}\label{lem:copies}
If $y^{(a)}_{cd}=\pi_{cd}(\eta)$ for all $a$ and all $(c,d)\in Q_a$, then
$\prod_a\chi_a(\eta,y^{(a)})=(-1)^{\sum_{a<c<b<d}\pi_{ab}\pi_{cd}}$.
\end{lemma}

\begin{proof}
$\sum_a\sum_{a<c<d}\pi_{cd}\sum_{c<b<d}\pi_{ab}=\sum_{a<c<b<d}\pi_{ab}\pi_{cd}$.
\end{proof}

So every interleaving pair of chords $\{a,b\}$, $\{c,d\}$ with $a<c<b<d$ is charged to node $a$,
which knows $\pi_{ab}$, and node $a$ receives a copy of $\pi_{cd}$ from node $c$, over the pair
$\{a,c\}$.

\subsection{Blob factorisation}
For $v\in V_a$ let $\deg_\eta(v):=\#\{k\in\partial(a): v\in\{s_k,t_k\},\ \eta_k\ne\lab N\}$, and
define the \emph{blob sums}
\[
  I_a(\eta):=\sum_{x\in\{0,1\}^{\Eint(a)}}\ \prod_{v\in V_a}\Ind{\deg_x(v)+\deg_\eta(v)=2}\cdot
  \topc{\omega_a}{\cF_a(x)\,R_a(\eta)},
\]
where
\[
  \cF_a(x):=\prod_{e\in\Eint(a)\colon x_e=1}(1+q_e)\in\Lambda_a\cap\Lambda^0 .
\]
$I_a$ depends only on $\eta|_{\partial(a)}$. An empty blob has $I_a=1$, and so does a blob that
consists of $r$ alone and has no edges.

\begin{theorem}[Blob factorisation]\label{thm:blob}
For every partition and every crossing order,
\[
  n\cdot\HC(G)=\sum_{\eta\in\cL^{\mx}}\varepsilon(\eta)\prod_{a=1}^{g}I_a(\eta)
  =\sum_{\eta\in\cL^{\mx}}\prod_{a=1}^{g}\sigma_a(\eta)\,\chi_a\bigl(\eta,\pi(\eta)\bigr)\,I_a(\eta),
\]
where $\chi_a(\eta,\pi(\eta))$ means $\chi_a(\eta,y^{(a)})$ with $y^{(a)}_{cd}=\pi_{cd}(\eta)$.
\end{theorem}

\begin{proof}
The second equality is \cref{prop:sign-factor} together with \cref{lem:copies}. For the first:
\begin{enumerate}[label=(\arabic*)]
\item Expanding $\prod_aI_a$ gives a sum over $x_{\mathrm{int}}\in\{0,1\}^{\bigcup_a\Eint(a)}$ of
the degree indicators times $\prod_a\topc{\omega_a}{\cF_aR_a}$.
\item By \cref{lem:block}, since $\cF_aR_a\in\Lambda_a$,
$\prod_a\topc{\omega_a}{\cF_aR_a}=\topc{\Omega_V}{\prod_a\cF_aR_a}$.
\item The $\cF_a$ are even and central, so $\prod_a(\cF_aR_a)=(\prod_a\cF_a)\cdot R_1\cdots R_g$.
\item By \cref{prop:sign-rule}, $\varepsilon(\eta)R_1\cdots R_g=\prod_k\alpha_{\eta_k}(k)
\beta_{\eta_k}(k)$.
\item Fix $x_{\mathrm{int}}$ and the support $T:=\{k:\eta_k\ne\lab N\}$. The degree indicators say
exactly that $x:=x_{\mathrm{int}}\cup\mathbf 1_T$ is a $2$-factor, and they do not depend on the
labels on $T$.
\item By distributivity and \cref{lem:pieces}(b),
$\sum_{\eta\colon\supp\eta=T}\prod_k\alpha_{\eta_k}(k)\beta_{\eta_k}(k)=\prod_{k\in T}(1+q_{e_k})$.
\item Hence $\sum_\eta\varepsilon\prod_aI_a=\sum_{x\in\{0,1\}^E}\Ind{x\text{ is a $2$-factor}}
\,\topc{\Omega_V}{\prod_{x_e=1}(1+q_e)}=n\cdot\HC(G)$ by \cref{cor:hc-berezin}.\qedhere
\end{enumerate}
\end{proof}

\begin{remark}\label{rem:sign-needed}
The interleaving sign matters exactly when odd labels sit on interleaving chords. In practice
this happens on disconnected $2$-factors whose non-root cycle winds through the blobs in
interleaving order; there the odd terms are what cancel $\kappa$ to $0$. Dropping it gives wrong
values on a positive fraction of test instances (\cref{tab:checks}, C12). Since every blob
carries an even number of odd generators, $(\pi_{ab})$ is a cycle-space vector of $K_g$, which
explains why the global sign is nontrivial only on a small fraction of random instances.
\end{remark}

% =====================================================================
\section{The network}\label{sec:network}
% =====================================================================

Fix $G$, $r$, $\phi$ and a crossing order as in \cref{sec:blobs}. Let
$\enc\colon\cL\to\{0,\dots,4\}$ be the encoding $\enc(\lab N)=0$ and $\enc(\lab j)=j+1$. For an
integer $t\ge0$ let $w(t):=\min\{w\ge0 : 2^w\ge 5^t\}=\lceil t\log 5\rceil$, and put
$w_{ab}:=w(m_{ab})$.

\begin{definition}[{The network $\mathcal N(G,r,\phi,\text{order})$ on $[g]$}]\label{def:our-network}
\leavevmode
\begin{description}[leftmargin=0pt,font=\normalfont\itshape]
\item[Settings.] A setting of node $a$ is a pair $z_a=(\eta^{(a)},y^{(a)})$ with
$\eta^{(a)}\in\cL^{\partial(a)}$ and $y^{(a)}\in\{0,1\}^{Q_a}$. It is encoded as the
concatenation of
\begin{itemize}
\item for each $b\in[g]\setminus\{a\}$, in increasing order, the \emph{block} $\block_b(z_a)$:
the $w_{ab}$-bit binary representation (lowest bit first) of
$\sum_{j=1}^{m_{ab}}\enc(\eta^{(a)}_{k_j})\,5^{j-1}$, where $k_1<\dots<k_{m_{ab}}$ enumerate
$E(\min(a,b),\max(a,b))$ in crossing order;
\item the bits $y^{(a)}_{cd}$ for $(c,d)\in Q_a$, in lexicographic order.
\end{itemize}
Thus $P_a:=\sum_{b\ne a}w_{ab}+|Q_a|$, and $S_a$ is the set of strings of length $P_a$ each of
whose blocks is $<5^{m_{ab}}$. These are exactly the encodings, and the encoding is a bijection,
so $|S_a|=5^{|\partial(a)|}2^{|Q_a|}$.
\item[Node functions.] $W_a(z_a):=\sigma_a(\eta^{(a)})\cdot\chi_a(\eta^{(a)},y^{(a)})\cdot
I_a(\eta^{(a)})$.
\item[Edges.] For $a<b$,
\[
  \Psi_{a,b}(z_a,z_b):=\prod_{k\in E(a,b)}\Ind{\eta^{(a)}_k=\eta^{(b)}_k}\cdot
  \prod_{d=b+1}^{g}\Ind{y^{(a)}_{bd}=\pi_{bd}(\eta^{(b)})},
\]
and $\Psi_{b,a}(z_b,z_a):=\Psi_{a,b}(z_a,z_b)$. The number of shared bits is
$c_{a,b}=c_{b,a}:=w_{ab}+(g-b)$, and the factorisation is $f_{a,b}(z_a,\theta):=
\Ind{h_{ab}(z_a)=\theta}$, $f_{b,a}(z_b,\theta):=\Ind{h_{ba}(z_b)=\theta}$, where
\begin{align*}
  h_{ab}(z_a)&:=\bigl(\block_b(z_a),\,y^{(a)}_{b,b+1},\dots,y^{(a)}_{b,g}\bigr)\in\{0,1\}^{c_{a,b}},\\
  h_{ba}(z_b)&:=\bigl(\block_a(z_b),\,\pi_{b,b+1}(\eta^{(b)}),\dots,\pi_{b,g}(\eta^{(b)})\bigr)\in\{0,1\}^{c_{a,b}}.
\end{align*}
\end{description}
\end{definition}

In words: the $w_{ab}$ block bits on $\{a,b\}$ are the base-$5$ code of the labels of the crossing
edges between $V_a$ and $V_b$; both ends hold a copy, and $\Psi$ tests that the copies agree. For
$a<b$ the remaining $g-b$ bits on $\{a,b\}$ are copies, held by $a$, of the parities $\pi_{bd}$
for $d>b$, which node $b$ computes from its own labels on $E(b,d)$.

\begin{lemma}\label{lem:factorisation}
$(f_{a,b},f_{b,a})$ factors $\Psi_{a,b}$ through $c_{a,b}$ bits.
\end{lemma}

\begin{proof}
Both ends list $E(a,b)$ in the same order, and the base-$5$ encoding is injective, so
$\block_b(z_a)=\block_a(z_b)$ iff $\eta^{(a)}$ and $\eta^{(b)}$ agree on $E(a,b)$. The remaining
coordinates of $h_{ab}$ and $h_{ba}$ agree iff every copy is correct. Hence
$\Psi_{a,b}=\Ind{h_{ab}(z_a)=h_{ba}(z_b)}$, and \cref{def:factor} applies.
\end{proof}

\begin{theorem}\label{thm:network-value}
The value of $\mathcal N(G,r,\phi,\text{order})$ is $n\cdot\HC(G)$, for every $n\ge3$, every
$\phi\colon V\to[g]$, every root $r$ and every crossing order.
\end{theorem}

\begin{proof}
The functions $\Psi_{a,b}$ are $0/1$-valued, so a term $\prod_aW_a\prod_{a<b}\Psi_{a,b}$ is
nonzero only if all $\Psi_{a,b}$ equal $1$. This means: (i) there is $\eta\in\cL^{\mx}$ with
$\eta^{(a)}=\eta|_{\partial(a)}$ for every $a$, because each $k$ lies in exactly $\partial(a_k)$
and $\partial(b_k)$, and $\Psi_{a_k,b_k}$ tests agreement on $k$; and (ii)
$y^{(a)}_{cd}=\pi_{cd}(\eta)$ for all $a<c<d$, because $\Psi_{a,c}$ tests this and
$\pi_{cd}(\eta^{(c)})=\pi_{cd}(\eta)$. Conversely, every $\eta\in\cL^{\mx}$ determines exactly one
$z$ with (i) and (ii). So $Z=\sum_\eta\prod_a\sigma_a(\eta)\chi_a(\eta,\pi(\eta))I_a(\eta)=
n\cdot\HC(G)$ by \cref{thm:blob}.
\end{proof}

\begin{proposition}[Parameters]\label{prop:parameters}
Let $\mu_a:=\sum_{v\in V_a}\deg_G(v)$ and $\mu:=\max_a\mu_a$, and suppose that $G$ has minimum
degree at least $2$. Then $\mathcal N(G,r,\phi,\text{order})$ is a network in the sense of
\cref{def:network} with
\[
  P:=w(\mu)+g+g^2,\qquad B:=2\mu+1,\qquad \tau:=2^{1.5\mu}n^{C_2}\quad(C_2\text{ absolute}),
\]
\[
  U=2\sum_{a<b}w_{ab}+2\binom g3\ \le\ 2\log5\cdot\mx+g^2+\frac{g^3}{3}\ \le\ 2\log5\cdot|E|+g^3,
\]
where $\tau$ is counted after a polynomial-time preparation. In particular $P\le\log5\cdot\mu+2g^2$
and $gP\le\log5\cdot g\mu+2g^3$.
\end{proposition}

\begin{proof}
\emph{Counting.} Each crossing edge at $a$ has exactly one end in $V_a$ and each internal edge
two, so $2|\Eint(a)|+|\partial(a)|=\mu_a$. Also $w(t)\le t\log5+\Ind{t>0}$ and $w(t)\ge t\log5$,
so $\sum_{b\ne a}w_{ab}\le|\partial(a)|\log5+(g-1)\le w(\mu)+g-1$.

\emph{$P_a\le P$.} $P_a=\sum_bw_{ab}+\binom{g-a}2\le w(\mu)+g-1+g^2/2\le P$.

\emph{Condition (i).} The copy bits on pairs at $a$ number
$\sum_{b>a}(g-b)+\sum_{b<a}(g-a)=\binom{g-a}2+(a-1)(g-a)\le g^2$, so $\sum_{b\ne a}c_{a,b}\le
w(\mu)+g-1+g^2\le P$.

\emph{The total.} $U=2\sum_{a<b}c_{a,b}$. Use $\sum_{a<b}(g-b)=\sum_b(b-1)(g-b)=\binom g3$ (choose
the middle element of a triple), and $\#\{a<b: m_{ab}>0\}\le g^2/2$. Finally $g^2+g^3/3\le g^3$ for
$g\ge2$.

\emph{Condition (ii).} $\Psi$ and $f$ take values in $\{0,1\}$, and $|\sigma_a\chi_a|=1$. In
$I_a$, the product $\cF_a(x)R_a(\eta)$ has at most $|\Eint(a)|$ factors with at most five signed
monomials each and $|\partial(a)|$ factors with at most two each, and there are at most
$2^{|\Eint(a)|}$ choices of $x$. So
\[
  |W_a|\le 10^{|\Eint(a)|}\,2^{|\partial(a)|}\le 2^{\frac{\log10}{2}\cdot2|\Eint(a)|+|\partial(a)|}
  \le 2^{1.67\mu_a}<2^B .
\]

\emph{Condition (iii).} The preparation stores the graph, $\phi$, the crossing order, the lists
$E(a,b)$, the numbers $w_{ab}$ and the sets $Q_a$. Then membership in $S_a$, encoding and
decoding, $h_{ab}$, $\Psi$ and $f$ take $\poly(n+g)$ operations (base-$5$ arithmetic on at most
$P$-bit numbers, with $P\le\poly(n)$); $\sigma_a$ takes $O(|\partial(a)|^2)$ operations and
$\chi_a$ takes $O(g^3)$. Finally, $I_a$ loops over at most $2^{|\Eint(a)|}\le2^{\mu/2}$ vectors
$x$; for each it multiplies at most $\mu$ factors (with at most five monomials each) in
$\Lambda_a$, represented densely by $2^{|\Gamma_a|}\le4^{|V_a|}\le2^\mu$ integers of at most
$B=O(n)$ bits (here $|V_a|\le\mu_a/2$ because the minimum degree is at least $2$), and reads one
coefficient. Altogether $\tau\le2^{\mu/2}\cdot2^\mu\cdot\poly(n)$.
\end{proof}

\begin{remark}\label{rem:largest-node}
\Cref{def:network} bounds the \emph{largest} node, through $P$, and not only the total. With
unbounded maximum degree one blob may have to carry $\Theta(n)$ crossing-edge bits, and then the
factor $2^{2gP/K}$ of \eqref{eq:lemma8} is not small. \Cref{sec:algorithm} removes high-degree
vertices first; the rank-encoded network of \cref{sec:degree} avoids the problem instead. Note also that the interleaving sign costs $O(g^3)$ bits in total and at most
$g^2$ bits per node.
\end{remark}

% =====================================================================
\section{The algorithm and its running time}\label{sec:algorithm}
% =====================================================================

\subsection{Preprocessing}

\begin{lemma}[Degree reduction]\label{lem:degree-reduction}
Let $\Theta\ge2$ and $H:=\{v:\deg(v)>\Theta\}$. For each choice $\nu=(\nu_v)_{v\in H}$ of a
$2$-element subset $\nu_v$ of the edges at $v$, let $G_\nu$ be $G$ minus every edge $e$ at a vertex
$v\in H$ with $e\notin\nu_v$. Then
\begin{enumerate}[label=\textup{(\alph*)}]
\item $\HC(G)=\sum_\nu\HC(G_\nu)$;
\item there are at most $\prod_{v\in H}\binom{\deg v}{2}\le n^{2|H|}$ choices, and $|H|<2m/\Theta$;
\item every $G_\nu$ has vertex set $V$, at most $m$ edges, and maximum degree at most $\Theta$.
\end{enumerate}
\end{lemma}

\begin{proof}
Map a Hamiltonian cycle $C$ of $G$ to $\nu(C):=(E(C)\cap\delta(v))_{v\in H}$, where $\delta(v)$ is
the set of edges at $v$. Then $C$ survives in $G_{\nu(C)}$, since every edge of $C$ at $v\in H$
lies in $\nu_v(C)$. Conversely, let $C$ be a Hamiltonian cycle of $G_\nu$. For $v\in H$, the set
$E(C)\cap\delta(v)$ has two elements and is contained in $\delta_{G_\nu}(v)\subseteq\nu_v$, so
$\nu=\nu(C)$. Thus the pairs $(\nu,C)$ with $C$ a Hamiltonian cycle of $G_\nu$ correspond
bijectively to the Hamiltonian cycles of $G$. The remaining claims hold because $\sum_v\deg
v=2m$ and the vertices of $H$ keep degree at most $2$.
\end{proof}

\begin{lemma}[Degree-balanced partition]\label{lem:balanced}
List $V$ as $v_1,\dots,v_n$ in any order, let $S_t:=\sum_{s<t}\deg(v_s)$, and put
$\phi(v_t):=\min\bigl(g,\lfloor gS_t/(2m)\rfloor+1\bigr)$. Then every blob has $\mu_a\le 2m/g+\Delta$,
where $\Delta$ is the maximum degree.
\end{lemma}

\begin{proof}
$\phi$ is nondecreasing in $t$, so each blob is a segment $v_t,\dots,v_{t'}$. If $a<g$, then
$(a-1)\cdot 2m/g\le S_t\le S_{t'}<a\cdot2m/g$, so $\mu_a=S_{t'}-S_t+\deg(v_{t'})<2m/g+\Delta$. If
$a=g$, then $\mu_g=2m-S_t\le 2m/g$.
\end{proof}

\subsection{The algorithm}
On input $G$ the algorithm does the following.
\begin{enumerate}[label=\textbf{Step \arabic*.},ref=\arabic*,leftmargin=*,start=0]
\item\label{st:0} If $n\le2$, or some vertex has degree less than $2$, output $0$.
\item\label{st:1} Choose a level $\ell$ (see \cref{sec:proof-delta,sec:proof-rate}). Put $K:=K_\ell$
and $g:=3K^8$.
\item\label{st:2} Apply \cref{lem:degree-reduction} with $\Theta:=\max(2,\lceil n/g\rceil)$. For each
branch $G_\nu$ do Steps~\ref{st:3}--\ref{st:4}; if $G_\nu$ has a vertex of degree less than $2$,
its count is $0$ and it is skipped.
\item\label{st:3} Take $\phi$ from \cref{lem:balanced}, any root $r$ and any crossing order, and
build $\mathcal N(G_\nu,r,\phi,\text{order})$ implicitly: perform the preparation of
\cref{prop:parameters}, computing $P$, $B$, $P_a$ and $c_{a,b}$ exactly.
\item\label{st:4} Run the procedure of \cref{imp:lemma8} on this network, with the algorithm of
\cref{imp:thmA}(b) at level $\ell$ as the clique algorithm; \cref{imp:lemma8} calls it on a
$(K,L)$-clique with $B'$ in the role of $B$. Divide the result by $n$.
\item\label{st:5} Output the sum over the branches.
\end{enumerate}
Correctness follows from \cref{lem:degree-reduction,thm:network-value,imp:lemma8,imp:thmA}.

\subsection{Checking the imported hypotheses}
\Cref{imp:lemma8} needs a network on $[g]$ with integer parameters $P,B,\tau\ge1$, total $U$, and
$g=3K^8$; \cref{prop:parameters} supplies these. It then needs an algorithm for every
$(K,L)$-clique with $|W|<2^{B'}$: \cref{imp:thmA}(b) at level $\ell$, with $\lambda=L$ and $B=B'$,
gives $T=2^{e_\ell K^2L}(K(L+B'))^{C_0K^2}$, and its hypotheses $\lambda,B\ge1$ hold because
$P\ge1$. \Cref{imp:thmA} needs pieces, supplied by \cref{imp:thmB}. Nothing in the clique
algorithm depends on the graph, which reaches it only through the tables that
\cref{imp:lemma8} builds from the $W_a$ and $f_{a,b}$. Balancing is done by \cref{imp:lemma8}
itself: it searches the $K^g$ groupings for one with at most $2g/K$ nodes per group and all
pairwise group totals at most $L$, which exists by the proof of \A[Lemma~1].

\subsection{Bounds for one branch}
Write $m:=|E(G)|$. Then $|E(G_\nu)|\le m$ and $G_\nu$ has minimum degree at least $2$. By
\cref{lem:balanced}, $\Delta\le\Theta\le n/g+2$, so $\mu\le(2m+n)/g+2$. Using $\log5\le2.33$ and
$g\ge3$, \cref{prop:parameters} gives
\begin{equation}\label{eq:branch-bounds}
\begin{gathered}
  gP\le2.33(2m+n)+3g^3,\qquad U\le4.65\,m+g^3,\qquad B\le 2(2m+n)/g+5,\\
  \log\tau\le1.5(2m+n)/g+3+C_2\log n,\qquad L\le gP+1,\qquad B'\le2gP+1+g^2B .
\end{gathered}
\end{equation}
By \eqref{eq:lemma8} and \eqref{eq:thmA}, with $e:=e_\ell\le1$ and $K^2L\le U+gP/K+K^2$, the time of
Steps~\ref{st:3}--\ref{st:4} for one branch is $2^X$ with
\begin{multline}\label{eq:exponent}
  X\le1+\max\Bigl\{g\log K,\ e\bigl(U+\tfrac{gP}{K}+K^2\bigr)+C_0K^2\log\bigl(K(L+B')\bigr)+\tfrac{2gP}{K}+\log\tau\Bigr\}\\
  +C_1\log\bigl(g(P+B)\bigr)+O(\log n).
\end{multline}
There are at most $n^{2|H|}\le n^{4mg/n}$ branches, because $|H|<2m/\Theta\le2mg/n$.

\subsection{Proof of \texorpdfstring{\cref{thm:main-delta}}{Theorem 1.1}}\label{sec:proof-delta}
Let $D\ge1$, $\delta\in(0,1]$ and $m\le Dn$, so that $2m+n\le3Dn$. By \cref{imp:thmA}(a),
$e_\ell\to0$ and $K_\ell\to\infty$, so there is $\ell$, depending only on $(\delta,D)$, with
$e:=e_\ell\le\delta/(40D)$ and $K:=K_\ell\ge60D/\delta$. Fix it, and put $g=3K^8$, a constant. By
\eqref{eq:branch-bounds}, the terms of \eqref{eq:exponent} satisfy
\begin{align*}
  e\,U&\le\tfrac{\delta}{40D}(4.65\,Dn+g^3)\le0.12\,\delta n+g^3,\\
  (e+2)\,gP/K&\le3\bigl(2.33\cdot3Dn+3g^3\bigr)/K\le21Dn/K+9g^3\le0.35\,\delta n+9g^3,\\
  \log\tau&\le4.5Dn/g+O(\log n)\le0.075\,\delta n+O(\log n),
\end{align*}
and the remaining terms are $O_{\delta,D}(\log n)$, since $g$ and $K$ are constants and $L$, $B'$,
$P$, $B$ are $O_{\delta,D}(n)$. So each branch costs $2^{0.55\delta n+O_{\delta,D}(\log n)}$. There
are at most $n^{4Dg}=n^{O_{\delta,D}(1)}$ branches, and the preparation is polynomial. The total
is $2^{0.55\delta n+O_{\delta,D}(\log n)}=O(2^{\delta n})$.\qed

The alternative route of \cref{rem:alt-route} gives the same bound: with $(K,e)$ achievable,
$e\le\delta/(40D)$ and $K\ge60D/\delta$, the exponent of $C_12^{e(U+gP/K)+2gP/K}(P+B)^{C_1}\tau$
is bounded by the same three lines.

\subsection{Proof of \texorpdfstring{\cref{thm:main-rate}}{Theorem 1.2}}\label{sec:proof-rate}
Use the pieces of \cref{imp:thmB} with $\rho=1/\sqrt2$, so $\alpha'=\frac12$, with $c$ and $x_0$
absolute. For $n<n_0:=K_0^{30}$, an absolute constant, any exact method takes $O(1)$ time.
Step~\ref{st:0} disposes of $n\le2$ and of minimum degree below $2$; otherwise $m\ge n$.

For $n\ge n_0$, let $\ell$ be the largest index with $K_\ell^{30}\le n$; $K_\ell$ is computable from
the formula in \cref{imp:thmA}(a). Then $K_{\ell+1}>n^{1/30}$, and $K_{\ell+1}\le K_\ell^{10}$ gives
$n^{1/300}<K:=K_\ell\le n^{1/30}$. So $g=3K^8\in[3n^{8/300},3n^{8/30}]$, and by \cref{imp:thmA}(a)
\[
  e=e_\ell\le0.85\Bigl(\frac{3x_0+6}{\log K}\Bigr)^{1/2}\le\frac{0.85\,(300(3x_0+6))^{1/2}}{\sqrt{\log n}}=:\frac{c_e}{\sqrt{\log n}} .
\]
With $2m+n\le3m$, the terms of \eqref{eq:exponent} and the branching satisfy
\begin{align*}
  e\,U&\le c_e\,(4.65\,m+27n^{0.8})/\sqrt{\log n},\\
  (e+2)\,gP/K&\le3\bigl(7m+81n^{0.8}\bigr)/n^{1/300}=O(m\,n^{-1/300}),\\
  \log\tau&\le4.5\,m/g+O(\log n)=O(m\,n^{-8/300}),\\
  eK^2+C_0K^2\log(K(L+B'))&=O(n^{1/15}\log n)\qquad(L,B'\le\poly(n),\text{ as }m\le n^2),\\
  g\log K&=O(n^{8/30}\log n),\qquad C_1\log(g(P+B))=O(\log n),\\
  \log(\#\text{branches})&\le4(m/n)\,g\log n=O(m\,n^{8/30-1}\log n).
\end{align*}
Every term except the first is $O(m/\sqrt{\log n})$ uniformly, because $m\ge n$ and $n^{0.8}$,
$n^{1/15}\log n$ and $n^{8/30}\log n$ are all $o(n/\sqrt{\log n})$. Hence the total time is
$2^{C^*m/\sqrt{\log n}}$ for an absolute $C^*$.\qed

\begin{remark}\label{rem:rate}
We record four comments on these bounds.
\begin{enumerate}[label=(\alph*)]
\item The requirement $g=3K^8\le n^{8/30}$ keeps the $g^3$ copy bits and the $n^{O(Dg)}$ branching
at $o(n/\sqrt{\log n})$; it forces only $K\ge n^{1/300}$, which still gives $e_\ell=
O(1/\sqrt{\log n})$.
\item With the pieces of \A[Lemma~2] only ($\rho=11/12$), the same argument gives
$2^{O(m/(\log n)^{0.125})}$.
\item The leading term is $e_\ell U\approx e_\ell\cdot4.64\,|E|$; every other term is of lower
order. Hence any improvement of $e_\ell$ in \cref{imp:thmA}, for example from better pieces,
transfers verbatim.
\item Like the bounds of \A, these are purely asymptotic. \Cref{thm:main-rate} beats $2^n$ only
when $e_\ell\cdot4.64\,m<n$, which needs $\log n$ in the thousands even for $m=2n$.
\end{enumerate}
\end{remark}

% =====================================================================
\section{Counting with logarithmic cost in the degree}\label{sec:degree}
% =====================================================================

This section proves \cref{thm:main-degree}. The proof is formalised in Lean (\lean{theorem3},
\lean{theorem3\_maxdeg}; files \lean{Unrank}, \lean{PairCode}, \lean{VertexBalanced},
\lean{RankNetwork}, \lean{RankProcs}, \lean{DegTableProc}, \lean{RankWordProcs}, \lean{Theorem3},
\lean{DPackage}). It reuses \cref{thm:blob} verbatim and changes only the network of
\cref{sec:network}: the pair $\{a,b\}$ no longer shares one base-$5$ digit for \emph{every}
crossing edge of $E(a,b)$, but only the set of crossing edges of $E(a,b)$ that carry a label other
than $\lab N$, rank-encoded among the subsets of the right size, together with the labels of those
edges. A $2$-factor uses exactly two of the $\deg(v)$ edges at each vertex $v$, so this costs about
$2\log\deg(v)$ bits per vertex rather than $\log_25$ bits per edge. The rank-encoding device is
that of \A[Appendix~D, Step~3]; we prove what we need from scratch. Fix $G$ with $n\ge3$ and
minimum degree at least $2$ (so $m\ge n$ and $\bar d\ge2$), a root $r$, a partition $\phi$ and a
crossing order, with the conventions of \cref{sec:blobs}. Put $n_a:=|V_a|$, $\bar n:=\max_an_a$,
$\lambda_\Delta:=\log(2e\Delta)$ and $\lambda_{\bar d}:=\log(2e\bar d)$; both are at least
$\log(4e)>3.4$.

\subsection{Profiles and the rank-encoded network}
A \emph{profile} is an integer vector $\xi=(\xi_{ab})_{a<b}$ with $0\le\xi_{ab}\le m_{ab}$; put
$\xi_{ba}:=\xi_{ab}$, $\xi_a:=\sum_{b\ne a}\xi_{ab}$ and $|\xi|:=\sum_{a<b}\xi_{ab}$. The profile of
a labelling $\eta\in\cL^{\mx}$ is $\xi(\eta)_{ab}:=\#\{k\in E(a,b):\eta_k\ne\lab N\}$. A profile is
\emph{admissible} if $\xi_a\le2n_a$ for every $a$. For $0\le s\le t$, the colex rank
$\{u_1<\dots<u_s\}\mapsto\sum_{i=1}^s\binom{u_i}{i}$ is a bijection from the $s$-subsets of
$\{0,\dots,t-1\}$ onto $\{0,\dots,\binom ts-1\}$, inverted greedily with at most $t$ binomial
comparisons; write $\mathrm{cw}(t,s):=\lceil\log\binom ts\rceil$.

\begin{definition}[{The network $\mathcal N_\xi$ on $[g]$}]\label{def:rank-network}
Fix a profile $\xi$. A setting of node $a$ is a pair $z_a=(\eta^{(a)},y^{(a)})$ with
$\eta^{(a)}\in\cL^{\partial(a)}$ such that $\#\{k\in E(a,b):\eta^{(a)}_k\ne\lab N\}=\xi_{ab}$ for every
$b\ne a$, and $y^{(a)}\in\{0,1\}^{Q_a}$. It is encoded as the concatenation of
\begin{itemize}
\item for each $b\ne a$, in increasing order, the block $\block_b(z_a)$ of width
$c'_{ab}:=\mathrm{cw}(m_{ab},\xi_{ab})+2\xi_{ab}$: if $k_1<\dots<k_{m_{ab}}$ enumerate the bundle of
$\{a,b\}$ in crossing order and $T_{ab}:=\{i:\eta^{(a)}_{k_i}\ne\lab N\}$, the block is the colex
rank of $T_{ab}$ in $\mathrm{cw}(m_{ab},\xi_{ab})$ bits, followed by the $2$-bit codes of the labels
$\eta^{(a)}_{k_i}\in\{\lab0,\lab1,\lab2,\lab3\}$, $i\in T_{ab}$ increasing;
\item the bits $y^{(a)}_{cd}$, $(c,d)\in Q_a$, in lexicographic order.
\end{itemize}
$S_a$ is the set of strings of length $P_a:=\sum_{b\ne a}c'_{ab}+|Q_a|$ whose rank field in every
block is $<\binom{m_{ab}}{\xi_{ab}}$. The functions $W_a=\sigma_a\chi_aI_a$ and $\Psi_{a,b}$ are those
of \cref{def:our-network}, and $c_{a,b}:=c'_{ab}+(g-\max(a,b))$, with $h_{ab},h_{ba}$ and $f_{a,b},
f_{b,a}$ as there, built from the new blocks.
\end{definition}

\begin{lemma}\label{lem:rank-encoding}
The encoding is a bijection from the settings of $a$ onto $S_a$, and $\block_b(z_a)=\block_a(z_b)$
if and only if $\eta^{(a)}$ and $\eta^{(b)}$ agree on the bundle of $\{a,b\}$. Hence
$\Psi_{a,b}=\Ind{h_{ab}(z_a)=h_{ba}(z_b)}$, and $(f_{a,b},f_{b,a})$ factors $\Psi_{a,b}$ through
$c_{a,b}$ bits.
\end{lemma}

\begin{proof}
Both ends list the bundle in the same order and use the same $\xi_{ab}$. The block determines
$T_{ab}$ and the labels on it, hence the restriction of $\eta^{(a)}$ to the bundle (the other labels
are $\lab N$); conversely a restriction with exactly $\xi_{ab}$ labels other than $\lab N$ gives exactly
one block, and every rank below $\binom{m_{ab}}{\xi_{ab}}$ is the rank of exactly one subset. The
rest is the proof of \cref{lem:factorisation}.
\end{proof}

\begin{theorem}\label{thm:rank-value}
$\sum_{\xi\text{ admissible}}\operatorname{value}(\mathcal N_\xi)=n\cdot\HC(G)$, and
$\operatorname{value}(\mathcal N_\xi)=0$ for every profile that is not admissible.
\end{theorem}

\begin{proof}
As in the proof of \cref{thm:network-value}, a nonzero term of $\operatorname{value}(\mathcal N_\xi)$
comes from a global $\eta\in\cL^{\mx}$ with correct copies, and by \cref{def:rank-network} that
$\eta$ has $\xi(\eta)=\xi$; conversely each such $\eta$ gives exactly one $z$. So
$\operatorname{value}(\mathcal N_\xi)=\sum_{\eta:\xi(\eta)=\xi}\prod_a\sigma_a\chi_a(\eta,\pi(\eta))I_a$.
The profiles partition $\cL^{\mx}$, so the sum over all profiles is $n\cdot\HC(G)$ by
\cref{thm:blob}. If $\xi_a>2n_a$, then every $\eta$ with $\xi(\eta)=\xi$ has
$\sum_{v\in V_a}\deg_\eta(v)=\xi_a>2n_a$, since each crossing edge at $a$ has exactly one end in
$V_a$; so some $v\in V_a$ has $\deg_\eta(v)\ge3$, every indicator in $I_a(\eta)$ vanishes, and the
value is $0$.
\end{proof}

\begin{proposition}[Parameters]\label{prop:rank-parameters}
Let $\xi$ be admissible. Then $\mathcal N_\xi$ is a network in the sense of \cref{def:network} with
\[
  P:=\lceil2\bar n\lambda_\Delta\rceil+g+g^2,\qquad B:=\lceil2\bar n\lambda_\Delta\rceil+1,\qquad
  \tau:=2^{2\bar n\lambda_\Delta}n^{C_2}\quad(C_2\text{ absolute}),
\]
$\tau$ counted after a polynomial-time preparation, and
\[
  U=2\sum_{a<b}c'_{ab}+2\binom g3\ \le\ 2n\log(2e\bar d)+g^3 .
\]
$P$ and $B$ do not depend on $\xi$.
\end{proposition}

\begin{proof}
We use $\binom NM\le(eN/M)^M$, Vandermonde's inequality
$\prod_i\binom{N_i}{M_i}\le\binom{\sum N_i}{\sum M_i}$, $\sum_i\lceil x_i\rceil\le\sum_ix_i+\#\{i:x_i>0\}$,
and the fact that $F_N(M):=M\log(eN/M)+2M$ is nondecreasing on $[0,N]$ (its derivative is
$\log(N/M)+2$).

\emph{Per node.} Let $N_a:=|\partial(a)|\le\Delta n_a$. Admissibility gives $\xi_a\le\min(2n_a,N_a)$,
so $\sum_{b\ne a}c'_{ab}\le\log\binom{N_a}{\xi_a}+2\xi_a+g-1\le F_{N_a}(\xi_a)+g-1$. If $2n_a\le N_a$,
then $F_{N_a}(\xi_a)\le F_{N_a}(2n_a)\le2n_a\log(e\Delta/2)+4n_a=2n_a\lambda_\Delta$; otherwise
$F_{N_a}(\xi_a)\le(\log e+2)N_a<2n_a\log(4e)\le2n_a\lambda_\Delta$. With the copy bits counted as in
\cref{prop:parameters}, $P_a\le P$ and $\sum_{b\ne a}c_{a,b}\le2n_a\lambda_\Delta+g-1+g^2\le P$.

\emph{Total.} $U=2\sum_{a<b}c'_{ab}+2\binom g3$ as in \cref{prop:parameters}, and
$\sum_{a<b}c'_{ab}\le\log\binom{\mx}{|\xi|}+2|\xi|+\binom g2$. Admissibility gives
$2|\xi|=\sum_a\xi_a\le2n$, and $|\xi|\le\mx\le m$. If $n\le\mx$, monotonicity gives
$\log\binom{\mx}{|\xi|}+2|\xi|\le F_{\mx}(n)\le n\log(em/n)+2n$; if $\mx<n$, it is at most
$(\log e+2)\mx<n\log(4e)\le n\log(4em/n)$, as $m\ge n$. Since $n\log(4em/n)=n\log(2e\bar d)$, we get
$U\le2n\log(2e\bar d)+g^2+2\binom g3\le2n\log(2e\bar d)+g^3$.

\emph{Condition (ii).} $\Psi$ and $f$ are $0/1$-valued and $|\sigma_a\chi_a|=1$. In $I_a(\eta)$ only
those $x$ contribute with $\deg_x(v)=2-\deg_\eta(v)$ for all $v\in V_a$. The map
$x\mapsto(x\cap\delta(v))_{v\in V_a}$ is injective, so there are at most
$\prod_{v\in V_a}\bigl(1+\deg v+\binom{\deg v}2\bigr)\le\Delta^{2n_a}$ of them (as $1+d+\binom d2\le d^2$
for $d\ge2$). For each, $|x|\le n_a$, and the top coefficient is at most the $\ell^1$-norm of
$\cF_a(x)R_a(\eta)$, which is at most $5^{|x|}2^{\xi_a}\le20^{n_a}$: the $\ell^1$-norm is
submultiplicative (structure constants $0,\pm1$), $1+q_e$ has norm at most $5$, and a piece has norm
at most $2$. Hence $\log|W_a|\le2n_a(\log\Delta+2.17)<B$.

\emph{Condition (iii).} The preparation stores the graph, $\phi$, the order, the bundles, the
binomials $\binom{m_{ab}}{\xi_{ab}}$ and the layout. Membership, encoding, decoding (one unranking per
block), $h_{ab}$, $\Psi$ and $f$ take $\poly(n)$ operations, $\sigma_a$ and $\chi_a$ as in
\cref{prop:parameters}. For $I_a$ the procedure enumerates the tuples $(x_v)_{v\in V_a}$, $x_v$ a
subset of size $2-\deg_\eta(v)$ of the internal edges at $v$, keeps the consistent ones (each chosen
edge chosen at both ends), and for each computes the top coefficient by choosing one signed
monomial from each factor, multiplying (sign by inversion count) and accumulating when the product
is $\pm\omega_a$. This takes at most $\Delta^{2n_a}\cdot20^{n_a}\cdot\poly(n)\le
2^{2n_a\lambda_\Delta}n^{C_2}$ operations, with a state of $\poly(n)$ bits.
\end{proof}

There are at most $\prod_{a<b}(\min(m_{ab},2\bar n)+1)\le(2\bar n+1)^{\binom g2}$ admissible profiles,
enumerated and tested in polynomial time each.

\subsection{The algorithm and the proof of \texorpdfstring{\cref{thm:main-degree}}{Theorem 1.3}}
\begin{enumerate}[label=\textbf{Step \arabic*.},leftmargin=*,start=0]
\item If $n\le2$, or some vertex has degree less than $2$, output $0$.
\item Choose the level $\ell$, $K:=K_\ell$ and $g:=3K^8$ as in \cref{sec:proof-rate}.
\item List $V$ as $v_1,\dots,v_n$ and put $\phi(v_t):=\lceil tg/n\rceil$, so that $n_a\le\lceil n/g\rceil$
(each blob is the set of integers in a half-open interval of length $n/g$). Take any root and any
crossing order, and do the preparation of \cref{prop:rank-parameters}.
\item For each admissible profile $\xi$, run the procedure of \cref{imp:lemma8} on $\mathcal N_\xi$
with the algorithm of \cref{imp:thmA}(b) at level $\ell$ as the clique algorithm.
\item Output the sum of the results divided by $n$.
\end{enumerate}
There is no degree reduction and no degree-balanced partition. Correctness follows from
\cref{thm:rank-value,imp:lemma8,imp:thmA}; the imported hypotheses are checked as in
\cref{sec:algorithm}, with \cref{prop:rank-parameters} in place of \cref{prop:parameters}, and
\cref{imp:lemma8} is applied to each $\mathcal N_\xi$ separately (only $U$ depends on $\xi$).

\begin{proof}[Proof of \cref{thm:main-degree}]
For $n<n_0$ use any exact method. For $n\ge n_0$, as in \cref{sec:proof-rate},
$n^{1/300}<K\le n^{1/30}$, $g\in[3n^{8/300},3n^{8/30}]$ and $e:=e_\ell\le c_e/\sqrt{\log n}$. By Step~2,
$\bar n\le n/g+1$, and $\lambda_\Delta\le\log(2en)\le2\log n$ because $\Delta<n$. By
\cref{prop:rank-parameters},
\[
  U\le2n\lambda_{\bar d}+g^3,\qquad gP\le2(n+g)\lambda_\Delta+3g^3,\qquad
  \log\tau\le2\bar n\lambda_\Delta+C_2\log n,
\]
and $L\le U+gP+1$, $B'\le2gP+1+g^2B$ are $\poly(n)$. The time for one profile is $2^X$ with $X$
bounded by \eqref{eq:exponent}, whose terms satisfy
\begin{align*}
  e\,U&\le c_e\,(2n\lambda_{\bar d}+27n^{0.8})/\sqrt{\log n},\\
  (e+2)\,gP/K&\le3\bigl(4(n+g)\log n+3g^3\bigr)/n^{1/300}=O(n^{1-1/300}\log n),\\
  \log\tau&\le4(n/g+1)\log n+C_2\log n=O(n^{1-8/300}\log n),\\
  eK^2+C_0K^2\log(K(L+B'))&=O(n^{1/15}\log n),\qquad g\log K=O(n^{8/30}\log n),
\end{align*}
and $C_1\log(g(P+B))=O(\log n)$. The number of admissible profiles is at most
$(2n/g+3)^{g^2/2}=2^{O(n^{16/30}\log n)}$. Every term except the first is $o(n/\sqrt{\log n})$ for
every $\Delta\le n$, and $\lambda_{\bar d}>3.4$, so the total time is
$2^{C^\sharp n\lambda_{\bar d}/\sqrt{\log n}}$ for an absolute $C^\sharp$. Finally $\bar d\le\Delta$.
\end{proof}

\begin{remark}\label{rem:degree}
\leavevmode
\begin{enumerate}[label=(\alph*)]
\item The degree enters in two ways. The total $U$ depends only on $\bar d$, for every partition.
The per-node parameters depend on $\Delta$, but only logarithmically, so a blob of $\lceil n/g\rceil$
vertices of degree $n-1$ still has only $O((n/g)\log n)$ bits, and the term $2gP/K$ of
\eqref{eq:lemma8}, which forced the degree reduction of \cref{sec:algorithm}
(\cref{rem:largest-node}), is $o(n/\sqrt{\log n})$ for every $\Delta$.
\item The leading term is $e_\ell U\approx2e_\ell n\log(2e\bar d)$, against
$e_\ell\cdot2\log_25\cdot m\approx2.32\,e_\ell\bar dn$ in \cref{thm:main-rate}; the new term is the
smaller one for $\bar d\ge4$. Better pieces in \A{} transfer verbatim.
\item The self-reductions of \cref{cor:decision-search} delete edges, which does not increase
$\bar d$; for paths, the added vertex raises $\bar d$ by at most $2$. So decision, search and paths
follow within the bound of \cref{thm:main-degree} times $n^2+1$; this extension is not machine-checked and is not claimed.
\item With a \emph{fixed} level, as in \cref{sec:proof-delta}, the same computation gives
$O(2^{\delta n})$ on graphs of maximum degree at most $\Delta_0$ as soon as
$e\le\delta/(8\log(2e\Delta_0))$ and $K\ge24\log(2e\Delta_0)/\delta$, which is weaker than the
requirement $e\le\delta/(40D)$ of \cref{thm:main-delta} when $D$ is large. The average-degree form
does not follow this way, because $gP/K\approx2n\lambda_\Delta/K$ is not $O(\delta n)$ when
$\Delta$ is unbounded. We do not use this.
\end{enumerate}
\end{remark}

% =====================================================================
\section{Decision, search and paths}\label{sec:corollaries}
% =====================================================================

\begin{proof}[Proof of \cref{cor:decision-search}]
\emph{Decision.} $G$ is Hamiltonian iff $\HC(G)>0$; one call of \cref{thm:main-delta} or
\cref{thm:main-rate} suffices.

\emph{Search, by self-reduction.} Let $H:=G$. For each edge $e$ of $G$ in turn, if
$\HC(H-e)>0$, set $H:=H-e$. The invariant is $\HC(H)>0$. Let $C$ be a Hamiltonian cycle of the
final $H$. Every intermediate $H$ contains the final one, hence $C$. Suppose some edge $f$ of the
final $H$ were not in $C$. When $f$ was examined, $H-f$ contained $C$, so $f$ would have been
deleted. Hence the final $H$ equals $C$. The machine runs over all $n^2$ vertex pairs (deleting a
non-edge changes nothing), so there are $n^2+1$ calls, each on a graph with $n$ vertices and at most
$m$ edges; the polynomial factor is absorbed into the constant.

\emph{Hamiltonian paths.} Let $G^+$ be $G$ plus a new vertex $w$ adjacent to all of $V$. For
$n\ge2$, the Hamiltonian paths of $G$ correspond bijectively to the Hamiltonian cycles of $G^+$,
so $\HP(G)=\HC(G^+)$. The graph $G^+$ has $n+1$ vertices and $m+n\le(D+1)(n+1)$ edges; the apex
$w$ has degree $n$, which \cref{thm:main-rate} handles directly; for the fixed-density form the
Lean proof applies \cref{thm:main-delta} to $G^+$ with $D+1$ in place of $D$. Search works as
above. (For $s$--$t$ paths one would join a new vertex to $s$ and $t$ only; that variant is not
machine-checked and not claimed.)
\end{proof}

% =====================================================================
\section{Results that rest only on the top-level theorems}\label{sec:parity}
% =====================================================================

This section concerns \cref{thm:parity,thm:randomized}. Unlike \cref{thm:main-delta}, they use the
labs' $\#$SAT theorems only as black boxes, through CNF encodings.

\subsection{The Cut\&Count identity}
Let $r\in V$. The CNF $\mathrm{cut}(G,r)$ has variables $x_e$ ($e\in E$) and $c_v$ ($v\in V$) and
imposes: at each vertex $v$, exactly two of the $x_e$ with $e\ni v$ are true; for each edge
$uw$, $x_{uw}\to(c_u\leftrightarrow c_w)$; and $c_r=\mathrm{false}$. Each local constraint becomes
clauses by forbidding its violating local patterns.

\begin{lemma}[Cut\&Count~\cite{CutAndCount11}]\label{lem:cutcount}
The number of models of $\mathrm{cut}(G,r)$ is $\sum_X2^{\cc(X)-1}$, the sum over the $2$-factors
$X$ of $G$, where $\cc(X)$ is the number of components of $(V,X)$. In particular it is congruent
to $\HC(G)$ modulo $2$.
\end{lemma}

\begin{proof}
A model is a pair $(X,c)$ where $X$ is a $2$-factor and $c\colon V\to\{0,1\}$ is constant on the
components of $(V,X)$ with $c(r)=0$. There are $2^{\cc(X)-1}$ such $c$, which is odd iff $X$ is
connected, i.e.\ a Hamiltonian cycle.
\end{proof}

For unbounded degree, the vertex constraint is replaced by a sequential counter (two bits per
prefix of the incident edges), giving a $5$-CNF on $5m+3n$ variables with the same model count.
The Lean development proves these identities (\lean{count_cutCNF},
\lean{count_cutCNF_mod_two}, \lean{count_bigCNF}).

\begin{proof}[Proof of \cref{thm:parity}]
Feed the $5$-CNF to the respective $\#$SAT algorithm. With \A, the formula has $N=5m+3n$
variables and at most $85m+11n+2$ clauses (after padding so that every variable occurs), and the
Main Theorem gives a count within $2^{C\cdot5m'/\sqrt{\log m'}}$ steps, where $N\le m'\le
85m+11n+2$ is the number of clauses. This is the Lean theorem
\lean{hamParity_allGraphs_anthropic}, whose hypothesis \lean{MainResults.MainTheorem} is
discharged by \A's Lean proof (\cref{sec:verification}). With~\cite{OpenAI26}, the Lean theorem
\lean{hamParity_allGraphs} gives one machine with time $C(L+1)^d2^{A\rho(5m+3n)}$ and space
polynomial in the encoding length $L$; its hypothesis \lean{exact_count_and_sat} is discharged
by OpenAI's Lean proof. The directed version uses one selected out-arc and in-arc per vertex
(\lean{dirHamParity_subexponential}). Building the CNF from $G$ takes polynomial time; that step
is not formalised as a machine.
\end{proof}

\subsection{Why colourings do not give exact counts}
Cut\&Count extracts $[X\text{ connected}]$ from a sum over consistent colourings with local
weights. The following barrier shows that such filters cannot be exact with few colours; this is
what motivates the Grassmann construction.

\begin{proposition}[machine-checked: \lean{colouring\_barrier\_graph}]\label{prop:colouring-barrier}
Work over a field of characteristic $0$. Suppose that weights $z_0,\dots,z_{q-1}$ on $q$ colours,
with the root coloured $0$ and every component $C\not\ni r$ of a $2$-factor contributing
$h(C)=\sum_jz_j^{|C|}$, give a sum equal to $\HC(G)$ up to a nonzero factor, for every graph $G$
on $n$ vertices. Then $q\ge n-4$.
\end{proposition}

\begin{proof}
We have $z_0\ne0$, since otherwise Hamiltonian cycles contribute $0$; normalise $z_0=1$. The root's
cycle has at least $3$ vertices, so a non-root cycle has length $k\in\{3,\dots,n-3\}$, and every
such length occurs: in the disjoint union of $C_{n-k}\ni r$ and $C_k$, the only $2$-factor is
disconnected. So the filter needs $\sum_{j=0}^{q-1}z_j^k=0$ for $k=3,\dots,n-3$. Drop the $z_j$
equal to $0$, and group the rest into distinct nonzero values $\lambda_i$ with positive integer
multiplicities. This gives an exponential polynomial with at most $q$ distinct nonzero bases and
positive coefficients that vanishes at $n-5$ consecutive integers. If $n-5\ge q$, the
Vandermonde matrix on those integers forces every coefficient to $0$, a contradiction because
$\lambda=1$ has multiplicity at least $1$.
\end{proof}

Conversely $q=n-2$ colours with $z_j=\zeta^j$, $\zeta$ a primitive $(n-2)$-th root of unity,
suffice. So exact counting by colourings needs about $\log n$ bits per vertex, and the exponent
becomes worse than brute force; the positive multiplicities vanish only in small characteristic,
which is how the modulo-$2$ filter escapes. The proposition does not cover vertex-dependent
weights.

\subsection{Randomized decision}

\begin{proof}[Proof of \cref{thm:randomized}\textup{(a)}]
Add to $\mathrm{cut}(G,r)$ an arbitrary CNF $\mathcal H$ on the edge variables. Then the model
count is congruent modulo $2$ to $\#\{X\in\mathrm{HC}(G):\mathcal H(\mathbf 1_X)\}$
(\lean{count_cutCNFWith_mod_two}); it is even if $G$ has no Hamiltonian cycle and odd if
exactly one Hamiltonian cycle satisfies $\mathcal H$. Take $\mathcal H$ to be a random sparse hash
in the sense of Calabro, Impagliazzo, Kabanets and Paturi~\cite{CIKP08}. Their isolation lemma~\cite[Lemma~3]{CIKP08} is
stated for solution sets of $k$-CNFs, but its proof uses only that the set is a set; the Lean
development re-proves it for arbitrary finite $S\subseteq\{0,1\}^N$ as a counting statement
(\lean{isolation_count}), together with the entropy estimate needed to make the isolating
fraction at least $\exp(-\delta N)$ for every $\delta>0$ (\lean{isolation_fraction}). With
$S=\mathrm{HC}(G)$ and $N=|E|\le\Delta n/2$ this gives soundness (\lean{decision_sound}), the
success fraction (\lean{hamDecision_fraction}) and one machine for all hashes
(\lean{hamDecision_subexponential}). Repeating $t\exp(\delta|E|)$ independent trials for each of
polynomially many parameter choices gives error at most $e^{-t}$; this outer loop is the only
part not modelled in Lean.
\end{proof}

\begin{remark}[The superseded part (b); not claimed]\label{rem:randomized-b}
Let $P_G(w)=\sum_{H\in\mathrm{HC}(G)}\prod_{e\in H}w_e$. It is multilinear with $0/1$
coefficients, so as a function on $\F_p^E$ it is nonzero iff $G$ is Hamiltonian. By the
Alon--Füredi theorem~\cite{AlonFuredi93}, a polynomial of degree $d$ that does not vanish on
$A_1\times\dots\times A_N$ is nonzero on at least
$\min\{\prod y_i : 1\le y_i\le|A_i|,\ \sum y_i\ge\sum|A_i|-d\}$ points; with $A_i=\F_p$ and $d=n$,
a uniform $w$ gives $P_G(w)\ne0$ with probability at least $p^{-n/(p-1)-1}$. The value
$P_G(w)\bmod p$ is the model count, modulo $p$, of a CNF with $p$ colours per vertex
($\lceil\log p\rceil$ bits with a comparator), the root coloured $0$, bitwise constraints
$x_e\to(c_u=c_v)$, the vertex constraint ``exactly two'', and weight bits $z_e$ with
$x_e=0\Rightarrow z_e=0$ and $x_e=1\Rightarrow z_e<w_e$; its count is
$\sum_Xp^{\cc(X)-1}\prod_{e\in X}w_e\equiv P_G(w)\pmod p$. It has $O(|E|\log p)$ clauses of width
$O(\log p)$. \A's Main Theorem, which is uniform in $k$, gives time $2^{O(|E|\log^2p/\sqrt{\log
|E|})}$ per test. Taking a prime $p\approx\sqrt{\log n}$ and $p^{n/(p-1)+1}\cdot t$ independent
tests gives one-sided error $e^{-t}$ in total time $2^{O_D(n(\log\log n)^2/\sqrt{\log n})}$.
This argument is not machine-checked (the Alon--Füredi theorem is, as far as we know, not in
Mathlib) and is recorded only for the history of the project.
\end{remark}

% =====================================================================
\section{Extensions (sketched)}\label{sec:extensions}
% =====================================================================

\emph{This section is not at the level of rigour of the rest of the paper.} The identities below
were checked exactly by brute force on small instances, and the blob networks were evaluated on
random partitions; the algorithmic consequences are sketched, and all of them are conditional on
$(\mathrm{H}_{\mathrm{A}})$ applied to general networks, exactly as for \cref{thm:main-delta}.

\subsection{The template}
All problems below fit one template:
\begin{equation}\label{eq:template}
  Z=\sum_{s,\xi}\prod_vC_v\bigl(s_v,\xi|_{\delta(v)}\bigr)\cdot
  \topc{\Omega}{\prod_vM_v(s_v)\prod_{e=uv}F_e(s_u,s_v,\xi_e)} ,
\end{equation}
where $s_v$ is a vertex state from a constant-size alphabet, $\xi_e$ an edge state ($\delta(v)$ is the set of edges at $v$), $C_v$ an
arbitrary local integer table, $M_v$ an even element in the generators of $v$ (for example the
mass term $1+w\bpsi_v\psi_v$, or $\bpsi_v\psi_v$, which deletes $v$), and $F_e$ an even element in
the generators of the ends of $e$, possibly involving pairwise scalars such as $z^{w_e}$. A root
has no generators. On a blob partition, each crossing edge is split by an exact rank
factorisation of its crossing matrix for each parity, a label being a term of this
factorisation; the edge costs $\log_2(\#\text{labels})$ bits. Because each $F_e$ is even, both
pieces of a term have the same parity, and the sign rule of \cref{sec:blobs} applies unchanged,
at a cost of $O(g^3)$ parity bits. Global counting constraints are handled by generating
variables, folded into one integer by Kronecker substitution; this changes only $B$, which stays
polynomial. Several connectivity predicates can be imposed at once by giving each its own
independent Grassmann field; the top coefficient then factors over the fields (sharing one field
would compute $\det(L_{X_1}+L_{X_2})$ instead).

\begin{table}[ht]
\centering
\small
\caption{Problems fitting the template \eqref{eq:template}. All identities were checked exactly on
small instances (brute force $=$ determinant identity $=$ blob network on random partitions with
$2$--$4$ blobs). ``Labels'' is the number of labels per crossing edge.}\label{tab:extensions}
\begin{tabular}{@{}>{\raggedright\arraybackslash}p{4.6cm}>{\raggedright\arraybackslash}p{9.6cm}c@{}}
\toprule
Problem & Identity (local constraints $\times$ one determinant) & Labels \\
\midrule
$\HC$ (this paper) & $n\HC=\sum_{\text{2-factors }X}\det L_X^{(r)}$ & $5$\\
directed $\HC$ & $\#\mathrm{DHC}=\sum_{\text{cycle covers }X}\det(D_{\mathrm{out}}-A)_X^{(r)}$ & $3$\\
TSP/ATSP, weights $\le\poly(n)$ & $n\sum_Hz^{w(H)}=\sum_Xz^{w(X)}\kappa(X)$ & $5$ / $3$\\
Hamiltonian $s$--$t$ paths & $\sum_{X:\ \deg s=\deg t=1,\text{ else }2}\det L_X^{(s)}$ & $5$\\
spanning trees with local degree constraints & $[z^{n-1}]\sum_Xz^{|X|}\prod_vc_v(\deg_Xv)\,\kappa(X)$ & $5$\\
Steiner tree (count, min.\ weight) & $\sum_{S\supseteq T}\det L_w(G[S])^{(t_0)}$ & $5$\\
induced forests, min.\ FVS (count) & $[\text{forest}]=2^{|E(S)|-|S|}[z^{|S|-|E(S)|}]\det(L_S+z\operatorname{diag}(2-\deg_S))$ & $5$\\
all forests, $T_G(2,1)$ and $T_G(x,1)$ & the same identity with $S=V$ and free edges & $4$\\
$k$-paths, $k$-cycles (all $k$) & coefficient extraction from $\det(L_X+wI)$ with markers & $5$\\
connected VC / connected DS & decision and minimum only; counts are $\kappa$-weighted & $4$ / $5$\\
\bottomrule
\end{tabular}
\end{table}

\noindent\emph{Machine-checked rows.} For directed Hamiltonian cycles and for ATSP the statement below
is proved in Lean, without the degree hypothesis and also in rate form for $\#\mathrm{DHC}$, by
splitting every vertex into an in-, a middle and an out-copy and applying
\cref{thm:main-delta,thm:main-rate,thm:tsp} to the resulting undirected graph (\lean{theoremDHC},
\lean{theoremDHC\_rate}, \lean{theoremATSP}, \lean{theoremATSP\_decode}; standard axioms only). The
directed identity of the table is formalised as well (\lean{card\_dirHamCycles\_eq\_sum},
\lean{dblob\_factorisation\_w}). For all other rows the proposition is a sketch and is not claimed.

\begin{proposition}[Sketch; not claimed except for the machine-checked rows]\label{prop:extensions}
Assume $(\mathrm{H}_{\mathrm{A}})$, and let $\mathcal P$ be one of the problems of
\cref{tab:extensions}, on graphs with $|E|\le Dn$ and either bounded maximum degree or the
hub-handling of \cref{rem:hubs}. Then for every $\delta>0$ the quantity stated in the table
(exact count, optimum, or $\kappa$-weighted count) can be computed deterministically in time
$2^{\delta n}$.
\end{proposition}

\begin{proof}[Proof sketch]
The identity in the table is an instance of \eqref{eq:template}; the blob network is built as in
\cref{sec:network}, with $U\approx 2\log_2(\#\text{labels})\cdot|E|+O(g^3)$ and $P\approx U/g+g^2$,
and the bookkeeping of \cref{sec:algorithm} applies verbatim. The identities have been checked
numerically; the uniform statements of the bit bounds, and the hub-handling, have not been written
out in full.
\end{proof}

\begin{remark}[Hubs]\label{rem:hubs}
A vertex of degree $\Theta(n)$ breaks the per-node bound $P$. For degree-$2$ problems (cycles,
paths, cycle covers) one can branch as in \cref{lem:degree-reduction}, or compress the hub's
realisable labellings into $O(\log d)$ bits per pair. For tree and forest problems one can split
the hub into a path of copies joined by \emph{forced} edges (factor $q_e$ instead of $1+q_e$),
which preserves the counts of subtrees and induced forests. A root vertex costs nothing, because
it has no generators. These devices were checked on small instances only.
\end{remark}

\subsection{The Tutte polynomial on the line \texorpdfstring{$y=1$}{y=1}}
On the line $y=1$, the Tutte polynomial is $T_G(x,1)=\sum_{F}(x-1)^{r(E)-|F|}$, the sum over the
forests $F\subseteq E$, where $r(E)=n-\cc(G)$; so counting forests by number of edges gives the
whole polynomial $T_G(x,1)$. The key identity is exact and short.

\begin{lemma}[Forest indicator; written proof, not machine-checked, not claimed]\label{lem:forest}
Let $H$ be a graph on a vertex set $S$ with edge set $X$, and let $\deg_H$ denote degrees in $H$.
Then
\[
  [z^{\,|S|-|X|}]\,\det\bigl(L_H+z\operatorname{diag}(2-\deg_H)\bigr)=
  \begin{cases}2^{|S|-|X|}&\text{if $H$ is a forest},\\ 0&\text{otherwise.}\end{cases}
\]
\end{lemma}

\begin{proof}
By the matrix-forest theorem, $\det(L_H+\operatorname{diag}(d))=\sum_{F}\prod_{T}\bigl(\sum_{v\in
T}d_v\bigr)$, the sum over the spanning forests $F$ of $H$ and the product over the trees $T$ of
$F$. With $d_v=z(2-\deg_Hv)$, a forest with $j$ trees contributes a multiple of $z^j$. Every
spanning forest has at least $\cc(H)\ge|S|-|X|$ trees, with equality throughout iff $H$ is itself a
forest and $F=H$. If $H$ is not a forest, the coefficient of $z^{|S|-|X|}$ is therefore $0$. If $H$
is a forest, the coefficient comes from $F=H$ alone and equals $\prod_T\sum_{v\in T}(2-\deg_Tv)=
\prod_T\bigl(2|T|-2(|T|-1)\bigr)=2^{\cc(H)}=2^{|S|-|X|}$.
\end{proof}

In Grassmann form the mass $-z\deg_H(v)\bpsi_v\psi_v$ is spread over the edges, so all factors are
local or pairwise and \eqref{eq:template} applies with four labels per crossing edge. Brand, Dell
and Roth~\cite{BrandDellRoth19} showed that, under $\#$ETH, for every real $x\ne1$ there is
$\varepsilon>0$ such that $T_G(x,1)$ has no $\exp(\varepsilon n)$-time algorithm on simple graphs
with $O(n)$ edges; this was the last open
line of the $\#$ETH dichotomy for the Tutte plane. \Cref{lem:forest} together with
\cref{prop:extensions} would give, conditional on $(\mathrm{H}_{\mathrm{A}})$ and with the caveats
of this section, a deterministic $2^{\delta n}$ algorithm for the whole polynomial $T_G(x,1)$ on
such graphs. (The published implication of course remains true; it is vacuous once $\#$ETH fails.)
The best previous algorithm computes the Tutte polynomial in vertex-exponential
time~\cite{BHKK08Tutte}.

\begin{remark}[Further problems]
A more general version of the template allows a constant number of global predicates (connected,
tree, path, single cycle, forest, minimum number of components, and directed analogues such as
branchings and strong connectivity), each certified by a determinant that is positive exactly when
the predicate holds. It covers, among others, connected feedback vertex set, connected odd cycle
transversal, minimum cycle cover, $k$-leaf out-branchings, maximum full-degree spanning trees and
graph-metric TSP. Exact counting is obtained when each certificate is constant on solutions; for
connected vertex covers, connected dominating sets and connected subgraphs, only
$\kappa$-weighted counts are obtained. All of this is exploratory and numerically checked on
instances with at most seven vertices.
\end{remark}

% =====================================================================
\section{Verification and formalisation status}\label{sec:verification}
% =====================================================================

We describe what has been checked, by what means, and what remains open. The aim is to make
clear how much trust each statement deserves.

\subsection{Machine-checked in Lean}
\begin{itemize}
\item \emph{Parity and randomized decision (\cref{thm:parity,thm:randomized}(a)).} The
Cut\&Count identities, the counter encoding for unbounded degree, the restriction by edge CNFs,
the Calabro--Impagliazzo--Kabanets--Paturi isolation lemma for arbitrary finite sets, its entropy
estimate, and the machine theorems are proved in Lean~4~\cite{Lean4} with Mathlib~\cite{Mathlib20},
taking the labs' top-level theorems as hypotheses.
\item \emph{Discharging the hypotheses.} OpenAI's development was built from source at a pinned
commit (Lean v4.34.1; 11{,}492 build jobs, no \lean{sorry}); \lean{#print axioms} reports only
\lean{propext}, \lean{Classical.choice} and \lean{Quot.sound} for \lean{exact_count_and_sat}.
The instantiated theorems (\lean{hamParity_allGraphs_unconditional},
\lean{hamDecision_unconditional}, \lean{dirHamParity_unconditional},
\lean{dirHamDecision_unconditional}, and others) are closed statements with the same three
axioms, and a kernel replay with \lean{leanchecker} of the full import closure succeeds.
Anthropic's development was likewise built from source (Lean v4.33.1; 9{,}768 jobs);
\lean{MainResults.mainTheorem} depends only on the three standard axioms, and so does
\lean{hamParity_allGraphs_anthropic} instantiated with it.
\item \emph{Exact-count formalisation.} The mathematical content of \cref{thm:blob} is
machine-checked. The Lean theorem \lean{blob_factorisation} states that, for every graph $G$, root
$r$, partition $\phi$ into blobs and crossing order,
$\sum_{\ell}\prod_a \sigma_a(\ell)\,\chi_a(\ell,P(\ell))\,I_a(\ell)=n\cdot\#\mathrm{HC}(G)$.
It rests only on \lean{propext}, \lean{Classical.choice} and \lean{Quot.sound}. Its proof
formalises:
\begin{itemize}
\item \cref{lem:2factor-trees};
\item the Grassmann algebra of \cref{sec:grassmann} (Koszul rule, block factorisation, linear
forms);
\item \cref{lem:berezin-cb} and \cref{lem:fermionic-mt}, including
$\det(\mathbf B_F)^2=[F\text{ is a spanning tree}]$;
\item \cref{prop:sign-rule};
\item the sign combinatorics (\cref{lem:disjoint-chords}, \cref{prop:sign-factor}, \cref{lem:copies});
\item the locality of the blob tables.
\end{itemize}
The network of \cref{sec:network} is also formalised in \A's own \lean{Net} format, inside \A's Lean
environment. Its value is proved to equal $n\cdot\#\mathrm{HC}(G)$ (\cref{thm:network-value}), and
its parameters satisfy \A's conditions (\cref{prop:parameters}, with $3$ bits per label instead of
$\log_2 5$). \A's procedure interface for networks is instantiated generically. The preprocessing
lemmas (\cref{lem:degree-reduction,lem:balanced}) are formalised. The streaming procedure for the blob tables is also formalised in \A's cost model. It is proved to
compute $\sigma_a\chi_aI_a$ within $2\cdot 8^{\mu_a}$ steps using states of linear size. On top of it,
the composition with \A's Lean theorems gives \cref{thm:main-delta} at the level of Turing machines
(\lean{theorem1_TM}): for every $\delta>0$ and $D$, one machine outputs $\#\mathrm{HC}(G)$ within
$C\cdot 2^{\delta n}$ steps whenever $|E|\le Dn$. The hypothesis bundling the network, its procedures and their bounds is then constructed from the
formalised network and procedures (a polynomial-time input reader, the layout, and the size bounds),
so \cref{thm:main-delta} is formally verified in Lean end to end (\lean{theorem1}). On input the
adjacency matrix of $G$, a Turing machine with finite alphabets outputs $\#\mathrm{HC}(G)$ in binary
within $\lfloor C\cdot 2^{\delta n}\rfloor$ steps whenever $|E|\le Dn$. The theorem rests only on
the three standard axioms and on \A's Lean development. The sharper single-algorithm bound of \cref{thm:main-rate} is formalised as well (\lean{theorem2}):
one machine, with the same axioms, outputs $\#\mathrm{HC}(G)$ within
$\lfloor C\cdot 2^{C(m+n)/\sqrt{\log_2(n+2)}}\rfloor$ steps for every graph $G$ with $m$ edges.
A fresh kernel replay with \lean{leanchecker} of the full import closure has passed for
both \lean{theorem1} and \lean{theorem2}.
\item \emph{Degree bound, corollaries, directed graphs.} \Cref{thm:main-degree} is \lean{theorem3}
and \lean{theorem3\_maxdeg}; \cref{cor:decision-search} is \lean{decision\_TM}, \lean{search\_TM},
\lean{hamPaths\_count\_TM} and the related machines; directed Hamiltonian cycles and ATSP are
\lean{theoremDHC}, \lean{theoremDHC\_rate}, \lean{theoremATSP}; \cref{prop:colouring-barrier} is
\lean{colouring\_barrier\_graph}. All use the standard axioms only.
\item \emph{Weighted layer (TSP).} \Cref{thm:tsp} is formally verified end to end
(\lean{theoremTSP}, \lean{theoremTSP\_decode}): for every $\delta>0$ and $D$ there are a Turing
machine $M$ and a constant $C$ such that, on the input word \lean{tspCode G W w} (the graph, $W$
in unary and the $n^2$ weights in binary), $M$ outputs the binary expansion of
$n\sum_H2^{b\,w(H)}$ within $\lfloor C\cdot2^{\delta n}(W+2)^C\rfloor$ steps whenever $|E|\le Dn$
and all weights are at most $W$; the decoding statements (digit $j$ equals $n$ times the number
of tours of weight $j$, the lowest non-zero digit is the optimum, its value over $n$ is the number
of optimal tours) are part of the formal theorem. The ingredients are the weighted blob
factorisation \lean{blob_factorisation_w}, the decoding lemma \lean{kronecker_decode}, the
weighted network \lean{hcNetW_kron} in \A's \lean{Net} format, and a weighted table procedure.
Standard axioms only; a fresh \lean{leanchecker} replay of its closure has passed.
\end{itemize}
Lean guarantees that the labs' proofs establish their formal statements; it does not guarantee
that those statements formalise the printed theorems faithfully, which is a matter for expert
review. \A{} itself notes that its Lean proof was compared with the printed one by AI models, not
mechanically.

\subsection{Numerical verification}
All identities of \cref{sec:grassmann,sec:trees,sec:blobs,sec:network} were implemented literally
and independently of the code used during discovery, in exact integer arithmetic: an exterior
algebra over $\Z$ with a \emph{randomised} generator order; the pieces, $\varepsilon$,
$\sigma_a$, $\pi_{ab}$ and $\chi_a$; the network of \cref{def:our-network} with fixed-length
bit-string settings; and two evaluators that use $W$, $\Psi$ and $f$ only as black boxes, one by
brute force over $\prod_aS_a$ and one by \A's grouping into $K=2$ or $3$ groups. Instances include
random graphs with arbitrary maps $\phi$ (empty blobs allowed), planted blob graphs, planted
Hamiltonian cycles, and instances with $g\in\{4,5\}$ whose planted Hamiltonian cycle, or whose
planted \emph{disconnected} $2$-factor, visits the blobs in interleaving order. Every instance uses
a random root, crossing order and generator order. \Cref{tab:checks} summarises one run (60
instances, $n\le12$); runs with several other seeds also pass.

\begin{table}[ht]
\centering
\small
\caption{Numerical checks (exact integers). Counts are for one run with 60 instances.}\label{tab:checks}
\begin{tabular}{@{}clr@{}}
\toprule
& Identity checked & Checks\\
\midrule
C1 & $n\HC=\sum_X\kappa(V,X)$ and $\kappa\in\{0,n\}$ (\cref{lem:2factor-trees}) & $60$ / $135$\\
C2 & $\kappa=\det L_X^{(r)}=\topc{\Omega_V}{\prod(1+q_e)}$ (\cref{lem:fermionic-mt}) & $135$\\
C3 & \cref{lem:berezin-cb} on random integer matrices & $200$\\
C4 & Koszul rule on random homogeneous sequences (\cref{lem:koszul}) & $300$\\
C5 & block factorisation (\cref{lem:block}) & $180$\\
C6 & \cref{lem:pieces}, both orientations and root ends & $5{,}216$\\
C7 & \cref{prop:sign-rule} as an equality of algebra elements & $162{,}875$\\
C8 & \cref{prop:sign-factor} and \cref{lem:copies} & $162{,}875$ each\\
C9 & \cref{prop:sign-rule} with fresh generators (no cancellation); \cref{lem:disjoint-chords} & $400$ / $1{,}033$\\
C10 & \cref{def:network} for the network: (i)--(iii), $U$, $|S_a|$, $|W_a|<2^B$, $\Psi=\sum_\theta ff'$ & $\approx594{,}000$\\
C11 & \cref{thm:blob}; network value by brute force; network value via $K$-grouping & $60$ / $60$ / $103$\\
C12 & ablations: no interleaving sign / no local signs / no sign: wrong on $20$ / $8$ / $18$ of $60$ & --\\
\bottomrule
\end{tabular}
\end{table}

A second script checks \cref{lem:degree-reduction} (on $109$ random graphs and several thresholds),
\cref{lem:balanced} in the exact form $g\mu_a\le2m+g\Delta$ (on $400$ graphs, $g\le12$), the
counting identities used in \cref{prop:parameters}, and the self-reduction of
\cref{cor:decision-search}. These checks test the identities and the network algebra, sign
included; they cannot test the asymptotic bookkeeping, which is checked by hand.

A third script checks the rank-encoded network of \cref{sec:degree} on top of the same library:
the colex rank, \cref{lem:rank-encoding} (exact count, injectivity, round trip, surjectivity on
random strings), the factorisation of $\Psi$, the bounds of \cref{prop:rank-parameters} per node and
in total, the profile count, \cref{thm:rank-value} (sum over admissible profiles equals
$n\cdot\HC(G)$ by brute force, evaluated through the shared bit strings only; inadmissible
profiles give $0$), and the degree-pruned computation of $I_a$ with its iteration counts. Besides
seeded random instances ($n\le9$, $g\le4$, including hubs of degree $n-1$), it runs four fixed
instances with $g=4$ and $n\in\{7,8,9\}$ on which the network without the interleaving sign is
wrong, and a built-in mutation test: with $\chi_a$ removed, or with the labels not shared, the
rank-encoded network must give a wrong value on at least one of them (it does on all four). Runs
with three seeds pass. An independent AI referee re-derived the construction and its parameter
bounds and found it sound; it also found that the first version of the script did not exercise the
interleaving sign, which led to the fixed instances and the mutation test.

\subsection{Independent referee}
An independent referee (an AI agent working from the one-paragraph claim) derived the sign rule
on its own and wrote its own code before reading the write-up. It checked about $500$ instances of
``Berezin integral $=$ determinant'', about $790$ blob evaluations and $297$ evaluations of a
literal \A-style network (with $g\le5$, random, very unequal and empty blobs, random roots, and in
half of the tests random integer edge weights), with no failures; removing the global sign gave
wrong answers on $8$ instances. Its verdict is that the construction and bookkeeping are sound,
conditional on \A{} (Theorem~A(b), Lemma~8(i) and the existence of pieces), and that the
unbounded-degree trap is handled correctly by \cref{lem:degree-reduction}. It raised one wording
point, about the label alphabet, which \cref{rem:five-labels} addresses. The referee did not build
or check the Lean files.

\subsection{What is pending}
\begin{itemize}
\item Human expert review of this paper and, above all, of the labs' manuscripts.
\item A human check that the formal statements of \lean{theorem1} and \lean{theorem2} say what
\cref{thm:main-delta,thm:main-rate} say (the definitions to read are listed in
\cref{sec:limitations}(b)); the proofs were written by AI agents and reviewed by an AI
coordinator.
\item A full write-up of the extensions of \cref{sec:extensions} other than TSP.
\end{itemize}

\subsection{Dependence on \texorpdfstring{\A}{[A]}}
If \A[Theorem~A(b)] or \A[Lemma~8] failed,
\cref{thm:main-delta,thm:main-rate,thm:main-degree,cor:decision-search} would lose their proofs, while \cref{sec:grassmann,sec:trees,sec:blobs,sec:network} would be
unaffected. The uniformity in \cref{thm:main-rate} uses the constant $c$ of \cref{imp:thmB}
through $x_0$; if $c$ is not explicit, the statement asserts the existence of one algorithm, as in
\A. The cost model is that of \A; our arithmetic on $O(n)$-bit integers costs $\poly(n)$
operations and is absorbed in $\tau$ and in the polynomial factors.

% =====================================================================
\section{Open problems}\label{sec:open}
% =====================================================================

\subsection{Dense graphs} The network of \cref{sec:network} has $U\approx2\log_25\cdot|E|$ bits,
and the rank-encoded network of \cref{sec:degree} has $U\approx2n\log(2e\bar d)$ bits, so
\cref{thm:main-degree} beats $2^n$ only for $\bar d=2^{o(\sqrt{\log n})}$, roughly. Nothing is
claimed for dense graphs beyond the classical $O^*(2^n)$~\cite{Karp82,Bax93} and its
$2^{n-\Omega(\sqrt{n/\log n})}$ refinement~\cite{Bjorklund16}. For $\bar d=\Theta(n)$, about $2n\log n$
bits is the information content of a Hamiltonian cycle, so for any network of this kind that
shares the cycle across the blobs the leading term $e_\ell U$ is $\Omega(e_\ell n\log n)$, which is
above $n$ unless $e_\ell=o(1/\log n)$. Is
there a subexponential algorithm for $\HC$ on dense graphs, under the same hypotheses?

\subsection{Connected subgraphs and \texorpdfstring{$T(1,2)$}{T(1,2)}} The determinant certificate is a weight, not a $0/1$
test, for objects that are connected but not trees. Exact counts of connected spanning subgraphs
($T_G(1,2)$), of connected vertex covers and of connected dominating sets, and the whole line
$x=1$ of the Tutte plane, are not obtained. A $\kappa$-free connectivity indicator that is local
in the sense of \eqref{eq:template} would settle these.

\subsection{Derandomization from the top-level theorems} \Cref{thm:main-delta} gives deterministic
decision, but it uses \A's internal lemmas. Can deterministic decision be obtained from the
top-level $\#$SAT theorems alone, which are machine-checked? Any such algorithm must use the
structure of the instance: if an algorithm accesses $S=\mathrm{HC}(G)\subseteq\{0,1\}^N$ only
through adaptive queries $\sum_{x\in S}w_i(x)\bmod q_i$ with $q_i\le Q$ and arbitrary weights
$w_i$, then a Davenport-constant argument shows that distinguishing $S=\emptyset$ from all nonempty
$S$ requires $2^N\le\operatorname{lcm}(1,\dots,Q)(1+F\ln Q)$ for $F$ queries, so $F\ge2^{N-o(N)}$
for $Q=2^{(\log N)^{1/2}}$. (This generalises the observation that parity queries alone need $2^N$
queries; it is related to the isolation bounds of~\cite{DKMW13}.) On the other hand, a fixed list
of $2^{o(n)}$ weight vectors works for all bounded-degree graphs on $n$ vertices, and hardness
assumptions for $\mathrm E$ against $\#\mathrm P$-oracle circuits would give deterministic
$2^{o(n)}$ decision via~\cite{ImpagliazzoWigderson97}, in the relativised form
of~\cite{KvM02}. So the question is one of explicit hitting
sets. These statements have written proofs but are not formalised.

\subsection{The exponent} The leading term of the exponent is $e_\ell U$, and $e_\ell=\Theta(\rho^\ell)$
is governed by the pieces of \A. Pieces with $\rho<1/\sqrt2$ would improve the exponent beyond
$1/\sqrt{\log n}$, and $\rho\to\frac12$ would give $2^{O(m/\log m)}$, and
$2^{O(n\log\bar d/\log n)}$ in \cref{thm:main-degree}. Can the $\sqrt{\log n}$ be
improved, for $\HC$ specifically or in general?

\subsection{Planar and minor-free graphs} A bound of $2^{\delta n}$ does not beat the known
$2^{O(\sqrt n)}$ on planar graphs. Does the blob network compose with balanced separators to give
$2^{o(\sqrt n)}$?

\subsection{Practicality} As in \A, the bounds are astronomically far from practical. A less
asymptotic analysis, or a different evaluation engine, would be needed for any concrete gain.

\section*{Acknowledgements}
This work was done with AI assistance (Claude Code, Anthropic): the construction, the proofs,
the numerical verification scripts, the independent referee pass and the Lean developments were
produced in sessions with Claude Code. The Lean components described in \cref{sec:verification}
are machine-checked; everything else has been checked numerically and by an independent AI
referee, but awaits human expert review. Responsibility for any errors rests with the author.

\bibliographystyle{amsplain}
\bibliography{refs}

\end{document}
